% Generated by tools/make_body.py; edit generator, not this file.
% UNPROOFREAD OCR TRANSCRIPTION: page-image verification remains required.
\chapter{Analytic Functions as Mappings}\label{chap:3}

% Original PDF physical page 65
\sourcepage{65}
A function $w = f(z)$ may be viewed as a mapping which represents a point z by its image w. The purpose of this chapter is to study, in a preliminary way, the special properties of mappings defined by analytic functions.

In order to carry out this program it is desirable to develop the underlying concepts with sufficient generality, for otherwise we would soon be forced to introduce a great number of ad hoc definitions whose mutual relationship would be far from clear. Since present-day students are exposed to abstraction and generality at quite an early stage, no apologies are needed. It is perhaps more appropriate to sound a warning that greatest possible generality should not become a purpose.

In the first section we develop the fundamentals of point set topology and metric spaces. There is no need to go very far, for our main concern is with the properties that are essential for the study of analytic functions. If the student feels that he is already thoroughly familiar with this material, he should read it only for terminology.

The author believes that proficiency in the study of analytic functions requires a mixture of geometric feeling and computational skill. The second and third sections, only loosely connected with the first, are expressly designed to develop geometric feeling by way of detailed study of elementary mappings. At the same time we try to stress rigor in geometric thinking, to the point where the geometric image becomes the guide but not the foundation of reasoning.


% Original PDF physical page 66
\sourcepage{66}
\setcounter{section}{0}\section{Elementary Point Set Topology}\index{Set|(}\label{sec:3.1}
The branch of mathematics which goes under the name of topology is concerned with all questions directly or indirectly related to continuity. The term is traditionally used in a very wide sense and without strict limits. Topological considerations are extremely important for the foundation of the study of analytic functions, and the first systematic study of topology was motivated by this need.

The logical foundations of set theory belong to another discipline. Our approach will be quite naive, in keeping with the fact that all our applications will be to very familiar objects. In this limited framework no logical paradoxes can occur.

\setcounter{section}{1}\setcounter{subsection}{0}\subsection{Sets and Elements}\label{sec:3.1.1}
In our language a set will be a collection of identifiable objects, its elements\index{Element}. The reader is familiar with the notation $x \in X$ which expresses that x is an element of X (as a rule we denote sets by capital letters and elements by small letters). Two sets are equal if and only if they have the same elements. X is a subset of Y if every element of X is also an element of Y, and this relationship is indicated by $X \subset Y$ or $Y \supset X$ (we do not exclude the possibility that X = Y). The empty set\index{Empty set}\index{Set!empty} is denoted by $\varnothing$.

A set can also be referred to as a space, and an element as a point\index{Point}. Subsets of a given space are usually called point sets. This lends a geometric flavor to the language, but should not be taken too literally. For instance, we shall have occasion to consider spaces whose elements are functions; in that case a “point” is a function.

The intersection\index{Intersection} of two sets X and Y, denoted by $X \cap Y$, is formed by all points which are elements of both X and Y. The union\index{Union} $X \cup Y$ consists of all points which are elements of either X or Y, including those which are elements of both. One can of course form the intersection and union of arbitrary collections of sets, whether finite or infinite in number.

The complement\index{Complement} of a set X consists of all points which are not in X; it will be denoted by $\sim X$. We note that the complement depends on the totality of points under consideration. For instance, a set of real numbers has one complement with respect to the real line and another with respect to the complex plane. More generally, if $X \subset Y$ we can consider the relative complement $Y \sim X$ which consists of all points that are in Y but not in X (we find it clearer to use this notation only when $X \subset Y$ ).

It is helpful to keep in mind the distributive laws
\begin{equation*}
X \cup (Y \cap Z) = (X \cup Y) \cap (X \cup Z)
\end{equation*}
\begin{equation*}
X \cap (Y \cup Z) = (X \cap Y) \cup (X \cap Z)
\end{equation*}
% Original PDF physical page 67
\index{Set|)}\sourcepage{67}
and the De Morgan laws\index{De Morgan laws}
\begin{equation*}
\begin{array}{l} \sim (X \cup Y) = \sim X \cap \sim Y \\ \sim (X \cap Y) = \sim X \cup \sim Y. \end{array}
\end{equation*}
These are purely logical identities, and they have obvious generalizations to arbitrary collections of sets.

\setcounter{section}{1}\setcounter{subsection}{1}\subsection{Metric Spaces}\label{sec:3.1.2}
For all considerations of limits and continuity it is essential to give a precise meaning to the terms “sufficiently near” and “arbitrarily near.” In the spaces R and C of real and complex numbers, respectively, such nearness can be expressed by a quantitative condition $|x - y| < \varepsilon$. For instance, to say that a set X contains all x sufficiently near to y means that there exists an $\varepsilon > 0$ such that $x \in X$ whenever $|x - y| < \varepsilon$. Similarly, X contains points arbitrarily near to y if to every $\varepsilon > 0$ there exists an $x \in X$ such that $|x - y| < \varepsilon$.

What we need to describe nearness in quantitative terms is obviously a distance $d(x,y)$ between any two points. We say that a set S is a metric space\index{Space!metric|(}\index{Metric space}\index{Space!metric} if there is defined, for every pair $x \in S$, $y \in S$, a nonnegative real number $d(x,y)$ in such a way that the following conditions are fulfilled:

1. $d(x,y) = 0$ if and only if $x = y$.

2. $d(y,x) = d(x,y)$.

3. $d(x,z)\leq d(x,y) + d(y,z).$

The last condition is the triangle inequality.

For instance, $\mathbf{R}$ and $\mathbf{C}$ are metric spaces with $d(x,y) = |x - y|$. The $n$-dimensional euclidean space $\mathbf{R}^n$ is the set of real $n$-tuples
\begin{equation*}
x = (x _ {1}, \dots , x _ {n})
\end{equation*}
with a distance defined by $d(x, y)^2 = \sum_{1}^{n} (x_i - y_i)^2$. We recall that we have defined a distance in the extended complex plane by
\begin{equation*}
d (z, z ^ {\prime}) = \frac {2 | z - z ^ {\prime} |}{\sqrt {(1 + | z | ^ {2}) (1 + | z ^ {\prime} | ^ {2})}}
\end{equation*}
(see Chap.~\hyperref[chap:1]{1}, Sec.~\hyperref[sec:1.2.4]{2.4}); since this represents the euclidean distance between the stereographic images on the Riemann sphere, the triangle inequality is obviously fulfilled. An example of a function space is given by $C[a,b]$, the set of all continuous functions defined on the interval $a \leq x \leq b$. It becomes a metric space if we define distance by $d(f,g) = \max |f(x) - g(x)|$.

In terms of distance, we introduce the following terminology: For any $\delta > 0$ and any $y \in S$, the set $B(y, \delta)$ of all $x \in S$ with $d(x, y) < \delta$ is called
% Original PDF physical page 68
\sourcepage{68}
the ball\index{Ball} with center y and radius $\delta$. It is also referred to as the $\delta$ -neighborhood\index{Neighborhood} of y. The general definition of neighborhood is as follows:

\begin{originalstatement}{Definition}{1}\label{def:3.1}
A set $N \subset S$ is called a neighborhood of $y \in S$ if it contains a ball $B(y, \delta)$.
\end{originalstatement}
In other words, a neighborhood of y is a set which contains all points sufficiently near to y. We use the notion of neighborhood to define open set\index{Open set}:

\begin{originalstatement}{Definition}{2}\label{def:3.2}
A set is open if it is a neighborhood of each of its elements.
\end{originalstatement}
The definition is interpreted to mean that the empty set is open (the condition is fulfilled because the set has no elements). The following is an immediate consequence of the triangle inequality:

Every ball is an open set.

Indeed, if $z \in B(y,\delta)$, then $\delta' = \delta - d(y,z) > 0$. The triangle inequality shows that $B(z,\delta') \subset B(y,\delta)$, for $d(x,z) < \delta'$ gives $d(x,y) < \delta' + d(y,z) = \delta$. Hence $B(y,\delta)$ is a neighborhood of $z$, and since $z$ was any point in $B(y,\delta)$ we conclude that $B(y,\delta)$ is an open set. For greater emphasis a ball is sometimes referred to as an open ball, to distinguish it from the closed ball\index{Ball!closed} formed by all $x \in S$ with $d(x,y) \leq \delta$.

In the complex plane $B(z_{0},\delta)$ is an open disk with center $z_{0}$ and radius $\delta$; it consists of all complex numbers z which satisfy the strict inequality $|z - z_{0}| < \delta$. We have just proved that it is an open set, and the reader is urged to interpret the proof in geometric terms.

The complement of an open set is said to be closed. In any metric space\index{Space!metric} the empty set and the whole space are at the same time open and closed, and there may be other sets with the same property.

The following properties of open and closed sets\index{Closed set}\index{Set!closed} are fundamental:

The intersection of a finite number of open sets is open.

The union of any collection of open sets is open.

The union of a finite number of closed sets is closed.

The intersection of any collection of closed sets is closed.

The proofs are so obvious that they can be left to the reader. It should be noted that the last two statements follow from the first two by use of the De Morgan laws.

There are many terms in common usage which are directly related to the idea of open sets. A complete list would be more confusing than helpful, and we shall limit ourselves to the following: interior\index{Interior}, closure\index{Closure},


% Original PDF physical page 69
\sourcepage{69}
boundary\index{Boundary}, exterior\index{Exterior}.

(i) The interior of a set X is the largest open set contained in X. It exists, for it may be characterized as the union of all open sets $\subset X$. It can also be described as the set of all points of which X is a neighborhood. We denote it by Int X.

(ii) The closure of X is the smallest closed set which contains X, or the intersection of all closed sets $\supset X$. A point belongs to the closure of X if and only if all its neighborhoods intersect X. The closure is usually denoted by $X^{-}$, infrequently by Cl X.

(iii) The boundary of X is the closure minus the interior. A point belongs to the boundary if and only if all its neighborhoods intersect both X and $\sim X$. Notation: Bd X or $\partial X$.

(iv) The exterior of X is the interior of $\sim X$. It is also the complement of the closure. As such it can be denoted by $\sim X^{-}$.

Observe that $\operatorname{Int} X \subset X \subset X^{-}$ and that $X$ is open if $\operatorname{Int} X = X$, closed if $X^{-} = X$. Also, $X \subset Y$ implies $\operatorname{Int} X \subset \operatorname{Int} Y$, $X^{-} \subset Y^{-}$. For added convenience we shall also introduce the notions of isolated point\index{Isolated point}\index{Point!isolated} and accumulation point\index{Accumulation point}\index{Point!accumulation}. We say that $x \in X$ is an isolated point of $X$ if $x$ has a neighborhood whose intersection with $X$ reduces to the point $x$. An accumulation point is a point of $X^{-}$ which is not an isolated point. It is clear that $x$ is an accumulation point of $X$ if and only if every neighborhood of $x$ contains infinitely many points from $X$.

\subsection*{EXERCISES}
\originalnumber{1}{exercise:3.1.2.1.1}
1. If S is a metric space\index{Space!metric} with distance function $d(x,y)$, show that S with the distance function $\delta(x,y)=d(x,y)/[1+d(x,y)]$ is also a metric space. The latter space is bounded in the sense that all distances lie under a fixed bound.

\originalnumber{2}{exercise:3.1.2.1.2}
2. Suppose that there are given two distance functions $d(x,y)$ and $d_{1}(x,y)$ on the same space S. They are said to be equivalent if they determine the same open sets. Show that d and $d_{1}$ are equivalent if to every $\varepsilon > 0$ there exists a $\delta > 0$ such that $d(x,y) < \delta$ implies $d_{1}(x,y) < \varepsilon$, and vice versa. Verify that this condition is fulfilled in the preceding exercise.

\originalnumber{3}{exercise:3.1.2.1.3}
3. Show by strict application of the definition that the closure of $|z - z_0| < \delta$ is $|z - z_0| \leq \delta$.

\originalnumber{4}{exercise:3.1.2.1.4}
4. If $X$ is the set of complex numbers whose real and imaginary parts are rational, what is Int $X$, $X^{-}$, $\partial X$?

\originalnumber{5}{exercise:3.1.2.1.5}
5. It is sometimes typographically simpler to write $X^{\prime}$ for $\sim X$. With this notation, how is $X^{\prime-\prime}$ related to $X$? Show that $X^{-\prime-\prime-\prime-\prime} = X^{-\prime-\prime}$.

\originalnumber{6}{exercise:3.1.2.1.6}
6. A set is said to be discrete if all its points are isolated. Show that a discrete set in $\mathbf{R}$ or $\mathbf{C}$ is countable.

\originalnumber{7}{exercise:3.1.2.1.7}
7. Show that the accumulation points of any set form a closed set.


% Original PDF physical page 70
\sourcepage{70}
\setcounter{section}{1}\setcounter{subsection}{2}\subsection{Connectedness}\label{sec:3.1.3}
If E is any nonempty subset of a metric space\index{Space!metric} S we may consider E as a metric space in its own right with the same distance function $d(x,y)$ as on all of S. Neighborhoods and open sets on E are defined as on any metric space, but an open set on E need not be open when regarded as a subset of S. To avoid confusion neighborhoods and open sets on E are often referred to as relative neighborhoods and relatively open sets. As an example, if we regard the closed interval $0 \leq x \leq 1$ as a subspace of R, then the semiclosed interval $0 \leq x < 1$ is relatively open, but not open in R. Henceforth, when we say that a subset E has some specific topological property, we shall always mean that it has this property as a subspace, and its subspace topology is called the relative topology.

Intuitively speaking, a space is connected if it consists of a single piece. This is meaningness unless we define the statement in terms of nearness. The easiest way is to give a negative characterization: S is not connected if there exists a partition $S = A \cup B$ into open subsets A and B. It is understood that A and B are disjoint and nonempty. The connectedness of a space is often used in the following manner: Suppose that we are able to construct two complementary open subsets A and B of S; if S is connected, we may conclude that either A or B is empty.

A subset $E \subset S$ is said to be connected if it is connected in the relative topology. At the risk of being pedantic we repeat:

\begin{originalstatement}{Definition}{3}\label{def:3.3}
A subset of a metric space is connected if it cannot be represented as the union of two disjoint relatively open sets none of which is empty.
\end{originalstatement}
If E is open, a subset of E is relatively open if and only if it is open. Similarly, if E is closed, relatively closed means the same as closed. We can therefore state: An open set is connected if it cannot be decomposed into two open sets, and a closed set is connected if it cannot be decomposed into two closed sets. Again, none of the sets is allowed to be empty.

Trivial examples of connected sets\index{Connected set|(}\index{Set!connected|(}\index{Connected set}\index{Set!connected} are the empty set and any set that consists of a single point.

In the case of the real line it is possible to name all connected sets. The most important result is that the whole line is connected, and this is indeed one of the fundamental properties of the real-number system.

An interval is defined by an inequality of one of the four types: $a < x < b$, $a \leq x < b$, $a < x \leq b$, $a \leq x \leq b$. For $a = -\infty$ or $b = +\infty$ this includes the semi-infinite intervals and the whole line.\hyperlink{fn:70:dagger}{\textsuperscript{\ensuremath{\dagger}}}

\begingroup\renewcommand{\thefootnote}{}\footnotetext{\hypertarget{fn:70:dagger}{}\textsuperscript{\ensuremath{\dagger}} We denote open intervals by $(a,b)$ and closed intervals by $[a,b]$. Another common practice is to denote open intervals by $]a,b[$ and semiclosed intervals by $]a,b]$ or $[a,b[$. It is always understood that $a<b$.}\endgroup

% Original PDF physical page 71
\index{Space!metric|)}\sourcepage{71}
\begin{originalstatement}{Theorem}{1}\index{Greatest lower bound (g.l.b.)|(}\index{Inf|(}\index{Least upper bound (l.u.b.)|(}\index{Sup|(}\label{thm:3.1}
The nonempty connected subsets of the real line are the intervals.
\end{originalstatement}
We reproduce one of the classical proofs, based on the fact that any monotone sequence has a finite or infinite limit.

Suppose that the real line $\mathbf{R}$ is represented as the union $\mathbf{R} = A \cup B$ of two disjoint closed sets. If neither is empty we can find $a_1 \in A$ and $b_1 \in B$; we may assume that $a_1 < b_1$. We bisect the interval $(a_1, b_1)$ and note that one of the two halves has its left end point in $A$ and its right end point in $B$. We denote this interval by $(a_2, b_2)$ and continue the process indefinitely. In this way we obtain a sequence of nested intervals $(a_n, b_n)$ with $a_n \in A$, $b_n \in B$. The sequences $\{a_n\}$ and $\{b_n\}$ have a common limit $c$. Since $A$ and $B$ are closed $c$ would have to be a common point of $A$ and $B$. This contradiction shows that either $A$ or $B$ is empty, and hence $\mathbf{R}$ is connected.

With minor modifications the same proof applies to any interval.

Before proving the converse we make an important remark. Let E be an arbitrary subset of R and call $\alpha$ a lower bound of E if $\alpha \leq x$ for all $x \in E$. Consider the set A of all lower bounds. It is evident that the complement of A is open. As to A itself it is easily seen that A is open whenever it does not contain any largest number. Because the line is connected, A and its complement cannot both be open unless one of them is empty. There are thus three possibilities: either A is empty, A contains a largest number, or A is the whole line. The largest number a of A, if it exists, is called the greatest lower bound of E; it is commonly denoted as g.l.b. x or inf x for $x \in E$. If A is empty, we agree to set $a = -\infty$, and if A is the whole line we set $a = +\infty$. With this convention every set of real numbers has a uniquely determined greatest lower bound; it is clear that $a = +\infty$ if and only if the set E is empty. The least upper bound, denoted as l.u.b. x or sup x for $x \in E$, is defined in a corresponding manner. \hyperlink{fn:71:dagger}{\textsuperscript{\ensuremath{\dagger}}}

Returning to the proof, we assume that E is a connected set\index{Connected set}\index{Set!connected} with the greatest lower bound a and the least upper bound b. All points of E lie between a and b, limits included. Suppose that a point $\xi$ from the open interval $(a,b)$ did not belong to E. Then the open sets defined by $x < \xi$ and $x > \xi$ cover E, and because E is connected, one of them must fail to meet E. Suppose, for instance, that no point of E lies to the left of $\xi$. Then $\xi$ would be a lower bound, in contradiction with the fact that a is the greatest lower bound. The opposite assumption would lead to a similar contradiction, and we conclude that $\xi$ must belong to E. It follows that E is an open, closed, or semiclosed interval with the end points a and b; the cases $a = -\infty$ and $b = +\infty$ are to be included.

\begingroup\renewcommand{\thefootnote}{}\footnotetext{\hypertarget{fn:71:dagger}{}\textsuperscript{\ensuremath{\dagger}} The supremum of a sequence was introduced already in Chap.~\hyperref[chap:2]{2}, Sec.~\hyperref[sec:2.2.1]{2.1}.}\endgroup

% Original PDF physical page 72
\index{Greatest lower bound (g.l.b.)|)}\index{Inf|)}\index{Least upper bound (l.u.b.)|)}\index{Sup|)}\sourcepage{72}
In\index{Bounded set|(}\index{Set!bounded|(} the course of the proof we have introduced the notions of greatest lower bound and least upper bound. If the set is closed and if the bounds are finite, they must belong to the set, in which case they are called the minimum\index{Minimum} and the maximum\index{Maximum}. In order to be sure that the bounds are finite we must know that the set is not empty and that there is some finite lower bound and some finite upper bound. In other words, the set must lie in a finite interval; such a set is said to be bounded. We have proved:

\begin{originalstatement}{Theorem}{2}\label{thm:3.2}
Any closed and bounded nonempty set of real numbers has a minimum and a maximum.
\end{originalstatement}
The structure of connected sets\index{Connected set}\index{Set!connected} in the plane is not nearly so simple as in the case of the line, but the following characterization of open connected sets contains essentially all the information we shall need.

\begin{originalstatement}{Theorem}{3}\label{thm:3.3}
A nonempty open set in the plane is connected if and only if any two of its points can be joined by a polygon which lies in the set.
\end{originalstatement}
The notion of a joining polygon is so simple that we need not give a formal definition.

We prove first that the condition is necessary. Let A be an open connected set, and choose a point $a \in A$. We denote by $A_{1}$ the subset of A whose points can be joined to a by polygons in A, and by $A_{2}$ the subset whose points cannot be so joined. Let us prove that $A_{1}$ and $A_{2}$ are both open. First, if $a_{1} \in A_{1}$ there exists a neighborhood $|z - a_{1}| < \varepsilon$ contained in A. All points in this neighborhood can be joined to $a_{1}$ by a line segment, and from there to a by a polygon. Hence the whole neighborhood is contained in $A_{1}$, and $A_{1}$ is open. Secondly, if $a_{2} \in A_{2}$, let $|z - a_{2}| < \varepsilon$ be a neighborhood contained in A. If a point in this neighborhood could be joined to a by a polygon, then $a_{2}$ could be joined to this point by a line segment, and from there to a. This is contrary to the definition of $A_{2}$, and we conclude that $A_{2}$ is open. Since A was connected either $A_{1}$ or $A_{2}$ must be empty. But $A_{1}$ contains the point a; hence $A_{2}$ is empty, and all points can be joined to a. Finally, any two points in A can be joined by way of a, and we have proved that the condition is necessary.

For future use we remark that it is even possible to join any two points by a polygon whose sides are parallel to the coordinate axes. The proof is the same.

In order to prove the sufficiency we assume that $A$ has a representation $A = A_1 \cup A_2$ as the union of two disjoint open sets. Choose $a_1 \in A_1$, $a_2 \in A_2$ and suppose that these points can be joined by a polygon in $A$.


% Original PDF physical page 73
\index{Bounded set|)}\index{Set!bounded|)}\sourcepage{73}
One of the sides of the polygon\index{Polygon} must then join a point in $A_{1}$ to a point in $A_{2}$, and for this reason it is sufficient to consider the case where $a_{1}$ and $a_{2}$ are joined by a line segment. This segment has a parametric representation $z = a_{1} + t(a_{2} - a_{1})$ where t runs through the interval $0 \leq t \leq 1$. The subsets of the interval 0 < t < 1 which correspond to points in $A_{1}$ and $A_{2}$, respectively, are evidently open, disjoint, and nonvoid. This contradicts the connectedness of the interval, and we have proved that the condition of the theorem is sufficient.

The theorem generalizes easily to $R^{n}$ and $C^{n}$.

\begin{originalstatement}{Definition}{4}\label{def:3.4}
A nonempty connected open set is called a region\index{Region}.
\end{originalstatement}
By Theorem~\hyperref[thm:3.3]{3} the whole plane, an open disk $|z - a| < \rho$, and a half plane are regions. The same is true of any $\delta$ -neighborhood in $R^{n}$. A region is the more-dimensional analogue of an open interval. The closure of a region is called a closed region\index{Closed region}\index{Region!closed}. It should be observed that different regions may have the same closure.

It happens frequently that we have to analyze the structure of sets which are defined very implicitly, for instance in the course of a proof. In such cases the first step is to decompose the set into its maximal connected components. As the name indicates, a component\index{Component} of a set is a connected subset which is not contained in any larger connected subset.

\begin{originalstatement}{Theorem}{4}\label{thm:3.4}
Every set has a unique decomposition into components.
\end{originalstatement}
If E is the given set, consider a point $a \in E$ and let $C(a)$ denote the union of all connected subsets of E that contain a. Then $C(a)$ is sure to contain a, for the set consisting of the single point a is connected. If we can show that $C(a)$ is connected, then it is a maximal connected set\index{Connected set}\index{Set!connected}, in other words a component. It would follow, moreover, that any two components are either disjoint or identical, which is precisely what we want to prove. Indeed, if $c \in C(a) \cap C(b)$, then $C(a) \subset C(c)$ by the definition of $C(c)$ and the connectedness of $C(a)$. Hence $a \in C(c)$, and by the same reasoning $C(c) \subset C(a)$, so that in fact $C(a) = C(c)$. Similarly $C(b) = C(c)$, and consequently $C(a) = C(b)$. We call $C(a)$ the component of a.

Suppose that $C(a)$ were not connected. Then we could find relatively open sets $A, B \neq \emptyset$ such that $C(a) = A \cup B$, $A \cap B = \emptyset$. We may assume that $a \in A$ while $B$ contains a point $b$. Since $b \in C(a)$ there is a connected set $E_0 \subset E$ which contains $a$ and $b$. The representation $E_0 = (E_0 \cap A) \cup (E_0 \cap B)$ would be a decomposition into relatively open subsets, and since $a \in E_0 \cap A$, $b \in E_0 \cap B$ neither part would be empty. This is a contradiction, and we conclude that $C(a)$ is connected.


% Original PDF physical page 74
\sourcepage{74}
\begin{originalstatement}{Theorem}{5}\index{Dense set|(}\index{Locally connected set|(}\label{thm:3.5}
In $\mathbf{R}^n$ the components of any open set are open.
\end{originalstatement}
This is a consequence of the fact that the $\delta$ -neighborhoods in $R^{n}$ are connected. Consider $a \in C(a) \subset E$. If E is open it contains $B(a, \delta)$ and because $B(a, \delta)$ is connected $B(a, \delta) \subset C(a)$. Hence $C(a)$ is open. A little more generally the assertion is true for any space S which is locally connected. By this we mean that any neighborhood of a point a contains a connected neighborhood of a. The proof is left to the reader.

In the case of $R^{n}$ we can conclude, furthermore, that the number of components is countable. To see this we observe that every open set must contain a point with rational coordinates. The set of points with rational coordinates is countable, and may thus be expressed as a sequence $\{p_{k}\}$. For each component $C(a)$, determine the smallest k such that $p_{k} \in C(a)$. To different components correspond different k. We conclude that the components are in one-to-one correspondence with a subset of the natural numbers, and consequently the set of components is countable.

For instance, every open subset of R is a countable union of disjoint open intervals.

Again, it is possible to analyze the proof and thereby arrive at a more general result. We shall say that a set E is dense in S if $E^{-} = S$, and we shall say that a metric space is separable if there exists a countable subset which is dense in S. We are led to the following result:

In a locally connected separable space\index{Space!separable} every open set is a countable union of disjoint regions.

\subsection*{EXERCISES}
\originalnumber{1}{exercise:3.1.3.2.1}
1. If $X \subset S$, show that the relatively open (closed) subsets of X are precisely those sets that can be expressed as the intersection of X with an open (closed) subset of S.

\originalnumber{2}{exercise:3.1.3.2.2}
2. Show that the union of two regions is a region if and only if they have a common point.

\originalnumber{3}{exercise:3.1.3.2.3}
3. Prove that the closure of a connected set\index{Connected set}\index{Set!connected} is connected.

\originalnumber{4}{exercise:3.1.3.2.4}
4. Let $A$ be the set of points $(x, y) \in \mathbf{R}^2$ with $x = 0$, $|y| \leq 1$, and let $B$ be the set with $x > 0$, $y = \sin 1/x$. Is $A \cup B$ connected?

\originalnumber{5}{exercise:3.1.3.2.5}
5. Let E be the set of points $(x,y) \in \mathbf{R}^{2}$ such that $0 \leq x \leq 1$ and either y = 0 or y = 1/n for some positive integer n. What are the components of E? Are they all closed? Are they relatively open? Verify that E is not locally connected.

\originalnumber{6}{exercise:3.1.3.2.6}
6. Prove that the components of a closed set are closed (use Ex.~\hyperref[exercise:3.1.3.2.3]{3}).

\originalnumber{7}{exercise:3.1.3.2.7}
7. A set is said to be discrete if all its points are isolated. Show that a discrete set\index{Discrete set}\index{Set!discrete} in a separable metric space is countable.


% Original PDF physical page 75
\index{Connected set|)}\index{Dense set|)}\index{Locally connected set|)}\index{Set!connected|)}\sourcepage{75}
\setcounter{section}{1}\setcounter{subsection}{3}\subsection{Compactness}\index{Set!compact|(}\label{sec:3.1.4}
The notions of convergent sequences and Cauchy sequences are obviously meaningful in any metric space. Indeed, we would say that $x_{n} \rightarrow x$ if $d(x_{n}, x) \rightarrow 0$, and we would say that $\{x_{n}\}$ is a Cauchy sequence if $d(x_{n}, x_{m}) \rightarrow 0$ as n and m tend to $\infty$. It is clear that every convergent sequence is a Cauchy sequence. For R and C we have proved the converse, namely that every Cauchy sequence is convergent (Chap.~\hyperref[chap:2]{2}, Sec.~\hyperref[sec:2.2.1]{2.1}), and it is not hard to see that this property carries over to any $R^{n}$. In view of its importance the property deserves a special name.

\begin{originalstatement}{Definition}{5}\label{def:3.5}
A metric space is said to be complete if every Cauchy sequence is convergent.
\end{originalstatement}
A subset is complete if it is complete when regarded as a subspace. The reader will find no difficulty in proving that a complete subset of a metric space is closed, and that a closed subset of a complete space\index{Space!complete} is complete.

We shall now introduce the stronger concept of compactness\index{Compactness|(}\index{Compactness}. It is stronger than completeness in the sense that every compact space or set is complete, but not conversely. As a matter of fact it will turn out that the compact subsets of R and C are the closed bounded sets. In view of this result it would be possible to dispense with the notion of compactness, at least for the purposes of this book, but this would be unwise, for it would mean shutting our eyes to the most striking property of bounded and closed sets of real or complex numbers. The outcome would be that we would have to repeat essentially the same proof in many different connections.

There are several equivalent characterizations of compactness, and it is a matter of taste which one to choose as definition. Whatever we do the uninitiated reader will feel somewhat bewildered, for he will not be able to discern the purpose of the definition. This is not surprising, for it took a whole generation of mathematicians to agree on the best approach. The consensus of present opinion is that it is best to focus the attention on the different ways in which a given set can be covered by open sets.

Let us say that a collection of open sets is an open covering\index{Open covering} of a set X if X is contained in the union of the open sets. A subcovering\index{Subcovering} is a subcollection with the same property, and a finite covering is one that consists of a finite number of sets. The definition of compactness reads:

\begin{originalstatement}{Definition}{6}\label{def:3.6}
A set X is compact if and only if every open covering of X contains a finite subcovering.
\end{originalstatement}
In this context we are thinking of $X$ as a subset of a metric space $S$,
% Original PDF physical page 76
\sourcepage{76}
and the covering is by open sets of S. But if U is an open set in S, then $U \cap X$ is an open subset of X (a relatively open set), and conversely every open subset of X can be expressed in this form (Sec.~\hyperref[sec:3.1.3]{1.3}, Ex. \hyperref[exercise:3.1.3.2.1]{1}). For this reason it makes no difference whether we formulate the definition for a full space or for a subset.

The property in the definition is frequently referred to as the Heine-Borel property\index{Heine-Borel property}. Its importance lies in the fact that many proofs become particularly simple when formulated in terms of open coverings.

We prove first that every compact space is complete. Suppose that $X$ is compact, and let $\{x_n\}$ be a Cauchy sequence in $X$. If $y$ is not the limit of $\{x_n\}$ there exists an $\varepsilon > 0$ such that $d(x_n,y) > 2\varepsilon$ for infinitely many $n$. Determine $n_0$ such that $d(x_m,x_n) < \varepsilon$ for $m,n \geq n_0$. We choose a fixed $n \geq n_0$ for which $d(x_n,y) > 2\varepsilon$. Then $d(x_m,y) \geq d(x_n,y) - d(x_m,x_n) > \varepsilon$ for all $m \geq n_0$. It follows that the $\varepsilon$-neighborhood $B(y,\varepsilon)$ contains only finitely many $x_n$ (better: contains $x_n$ only for finitely many $n$).

Consider now the collection of all open sets U which contain only finitely many $x_{n}$. If $\{x_{n}\}$ is not convergent, it follows by the preceding reasoning that this collection is an open covering of X. Therefore it must contain a finite subcovering, formed by $U_{1}, \ldots, U_{N}$. But that is clearly impossible, for since each $U_{i}$ contains only finitely many $x_{n}$ it would follow that the given sequence is finite.

Secondly, a compact set\index{Set!compact} is necessarily bounded (a metric space is bounded if all distances lie under a finite bound). To see this, choose a point $x_0$ and consider all balls $B(x_0, r)$. They form an open covering of $X$, and if $X$ is compact, it contains a finite subcovering; in other words, $X \subset B(x_0, r_1) \cup \dots \cup B(x_0, r_m)$, which means the same as $X \subset B(x_0, r)$ with $r = \max(r_1, \ldots, r_m)$. For any $x, y \in X$ it follows that $d(x, y) \leq d(x, x_0) + d(y, x_0) < 2r$, and we have proved that $X$ is bounded.

But boundedness is not all we can prove. It is convenient to define a stronger property called total boundedness:

\begin{originalstatement}{Definition}{7}\label{def:3.7}
A set X is totally bounded if, for every $\varepsilon > 0$, X can be covered by finitely many balls of radius $\varepsilon$.
\end{originalstatement}
This is certainly true of any compact set. For the collection of all balls of radius $\varepsilon$ is an open covering, and the compactness\index{Compactness} implies that we can select finitely many that cover X. We observe that a totally bounded set is necessarily bounded, for if $X \subset B(x_{1}, \varepsilon) \cup \cdots \cup B(x_{m}, \varepsilon)$, then any two points of X have a distance $<2\varepsilon + \max d(x_{i}, x_{j})$. (The preceding proof that any compact set is bounded becomes redundant.)

We have already proved one part of the following theorem:

\begin{originalstatement}{Theorem}{6}\label{thm:3.6}
A set is compact if and only if it is complete and totally bounded.
\end{originalstatement}

% Original PDF physical page 77
\sourcepage{77}
To prove the other part, assume that the metric space S is complete and totally bounded. Suppose that there exists an open covering which does not contain any finite subcovering. Write $\varepsilon_{n}=2^{-n}$. We know that S can be covered by finitely many $B(x,\varepsilon_{1})$. If each had a finite subcovering, the same would be true of S; hence there exists a $B(x_{1},\varepsilon_{1})$ which does not admit a finite subcovering. Because $B(x_{1},\varepsilon_{1})$ is itself totally bounded we can find an $x_{2}\in B(x_{1},\varepsilon_{1})$ such that $B(x_{2},\varepsilon_{2})$ has no finite subcovering. \hyperlink{fn:77:dagger}{\textsuperscript{\ensuremath{\dagger}}} It is clear how to continue the construction: we obtain a sequence $x_{n}$ with the property that $B(x_{n},\varepsilon_{n})$ has no finite subcovering and $x_{n+1}\in B(x_{n},\varepsilon_{n})$. The second property implies $d(x_{n},x_{n+1})<\varepsilon_{n}$ and hence $d(x_{n},x_{n+p})<\varepsilon_{n}+\varepsilon_{n+1}+\cdots+\varepsilon_{n+p-1}<2^{-n+1}$. It follows that $x_{n}$ is a Cauchy sequence. It converges to a limit y, and this y belongs to one of the open sets U in the given covering. Because U is open, it contains a ball $B(y,\delta)$. Choose n so large that $d(x_{n},y)<\delta/2$ and $\varepsilon_{n}<\delta/2$. Then $B(x_{n},\varepsilon_{n})\subset B(y,\delta)$, for $d(x,x_{n})<\varepsilon_{n}$ implies $d(x,y)\leq d(x,x_{n})+d(x_{n},y)<\delta$. Therefore $B(x_{n},\varepsilon_{n})$ admits a finite subcovering, namely by the single set U. This is a contradiction, and we conclude that S has the Heine-Borel property.

\begin{originalstatement}{Corollary}{}\label{cor:3.unnumbered-p77-b1}
A subset of $\mathbf{R}$ or $\mathbf{C}$ is compact if and only if it is closed and bounded.
\end{originalstatement}
We have already mentioned this particular consequence. In one direction the conclusion is immediate: We know that a compact set\index{Set!compact} is bounded and complete; but R and C are complete, and complete subsets of a complete space are closed. For the opposite conclusion we need to show that every bounded set in R or C is totally bounded. Let us take the case of C. If X is bounded it is contained in a disk, and hence in a square. The square can be subdivided into a finite number of squares with arbitrarily small side, and the squares can in turn be covered by disks with arbitrarily small radius. This proves that X is totally bounded, except for a small point that should not be glossed over. When Definition~\hyperref[def:3.7]{7} is applied to a subset $X \subset S$ it is slightly ambiguous, for it is not clear whether the $\varepsilon$ -neighborhoods should be with respect to X or with respect to S; that is, it is not clear whether we require their centers to lie on X. It happens that this is of no avail. In fact, suppose that we have covered X by $\varepsilon$ -neighborhoods whose centers do not necessarily lie on X. If such a neighborhood does not meet X it is superfluous, and can be dropped. If it does contain a point from X, then we can replace it by a 2 $\varepsilon$ -neighborhood around that point, and we obtain a finite covering by 2 $\varepsilon$ -neighborhoods with centers on X. For this reason the ambiguity is only apparent, and our proof that bounded subsets of C are totally bounded is valid.

\begingroup\renewcommand{\thefootnote}{}\footnotetext{\hypertarget{fn:77:dagger}{}\textsuperscript{\ensuremath{\dagger}} Here we are using the fact that any subset of a totally bounded set\index{Set!totally bounded}\index{Totally bounded set} is totally bounded. The reader should prove this.}\endgroup

% Original PDF physical page 78
\sourcepage{78}
There is a third characterization of compact sets\index{Set!compact}. It deals with the notion of limit point\index{Limit point}\index{Point!limit} (sometimes called cluster value): We say that y is a limit point of the sequence $\{x_{n}\}$ if there exists a subsequence $\{x_{n_{k}}\}$ that converges to y. A limit point is almost the same as an accumulation point of the set formed by the points $x_{n}$, except that a sequence permits repetitions of the same point. If y is a limit point, every neighborhood of y contains infinitely many $x_{n}$. The converse is also true. Indeed, suppose that $\varepsilon_{k} \to 0$. If every $B(y, \varepsilon_{k})$ contains infinitely many $x_{n}$ we can choose subscripts $n_{k}$, by induction, in such a way that $x_{n_{k}} \in B(y, \varepsilon_{k})$ and $n_{k+1} > n_{k}$. It is clear that $\{x_{n_{k}}\}$ converges to y.

\begin{originalstatement}{Theorem}{7}\label{thm:3.7}
A metric space is compact if and only if every infinite sequence has a limit point.
\end{originalstatement}
This theorem is usually referred to as the Bolzano-Weierstrass theorem\index{Bolzano-Weierstrass theorem}. The original formulation was that every bounded sequence of complex numbers has a convergent subsequence. It came to be recognized as an important theorem precisely because of the role it plays in the theory of analytic functions.

The first part of the proof is a repetition of an earlier argument. If y is not a limit point of $\{x_{n}\}$ it has a neighborhood which contains only finitely many $x_{n}$ (abbreviated version of the correct phrase). If there were no limit points the open sets containing only finitely many $x_{n}$ would form an open covering. In the compact case we could select a finite subcovering, and it would follow that the sequence is finite. The previous time we used this reasoning was to prove that a compact space is complete. We showed in essence that every sequence has a limit point, and then we observed that a Cauchy sequence with a limit point is necessarily convergent. For strict economy of thought it would thus have been better to prove Theorem~\hyperref[thm:3.7]{7} before Theorem~\hyperref[thm:3.6]{6}, but we preferred to emphasize the importance of total boundedness as early as possible.

It remains to prove the converse. In the first place it is clear that the Bolzano-Weierstrass property implies completeness. Indeed, we just pointed out that a Cauchy sequence with a limit point must be convergent. Suppose now that the space is not totally bounded. Then there exists an $\varepsilon > 0$ such that the space cannot be covered by finitely many $\varepsilon$-neighborhoods. We construct a sequence $\{x_n\}$ as follows: $x_1$ is arbitrary, and when $x_1, \ldots, x_n$ have been selected we choose $x_{n+1}$ so that it does not lie in $B(x_1, \varepsilon) \cup \cdots \cup B(x_n, \varepsilon)$. This is always possible because these neighborhoods do not cover the whole space. But it is clear that $\{x_n\}$ has no convergent subsequence, for $d(x_m, x_n) > \varepsilon$ for all $m$ and $n$. We conclude that the Bolzano-Weierstrass property implies total boundedness. In view of Theorem~\hyperref[thm:3.6]{6} that is what we had to prove.


% Original PDF physical page 79
\sourcepage{79}
The\index{Into, definition|(}\index{Weierstrass K@Weierstrass, K.|(} reader should reflect on the fact that we have exhibited three characterizations of compactness\index{Compactness} whose logical equivalence is not at all trivial. It should be clear that results of this kind are particularly valuable for the purpose of presenting proofs as concisely as possible.

\subsection*{EXERCISES}
\originalnumber{1}{exercise:3.1.4.3.1}
1. Give an alternate proof of the fact that every bounded sequence of complex numbers has a convergent subsequence (for instance by use of the limes inferior).

\originalnumber{2}{exercise:3.1.4.3.2}
2. Show that the Heine-Borel property can also be expressed in the following manner: Every collection of closed sets with an empty intersection contains a finite subcollection with empty intersection.

\originalnumber{3}{exercise:3.1.4.3.3}
3. Use compactness to prove that a closed bounded set of real numbers has a maximum.

\originalnumber{4}{exercise:3.1.4.3.4}
4. If $E_1 \supset E_2 \supset E_3 \supset \cdots$ is a decreasing sequence of nonempty compact sets\index{Set!compact}, then the intersection $\bigcap_{1}^{\infty} E_n$ is not empty (Cantor\index{Cantor G@Cantor, G.}'s lemma). Show by example that this need not be true if the sets are merely closed.

\originalnumber{5}{exercise:3.1.4.3.5}
5. Let $S$ be the set of all sequences $x = \{x_n\}$ of real numbers such that only a finite number of the $x_n$ are $\neq 0$. Define $d(x, y) = \max |x_n - y_n|$. Is the space complete? Show that the $\delta$-neighborhoods are not totally bounded.

\setcounter{section}{1}\setcounter{subsection}{4}\subsection{Continuous Functions}\label{sec:3.1.5}
We shall consider functions $f$ which are defined on a metric space $S$ and have values in another metric space $S'$. Functions are also referred to as mappings: we say that $f$ maps $S$ into $S'$, and we write $f: S \to S'$. Naturally, we shall be mainly concerned with real- or complex-valued functions; occasionally the latter are allowed to take values in the extended complex plane, ordinary distance being replaced by distance on the Riemann sphere.

The space S is the domain\index{Domain} of the function. We are of course free to consider functions f whose domain is only a subset of S, in which case the domain is regarded as a subspace. In most cases it is safe to slur over the distinction: a function on S and its restriction to a subset are usually denoted by the same symbol. If $X \subset S$ the set of all values $f(x)$ for $x \in S$ is called the image\index{Image} of X under f, and it is denoted by $f(X)$. The inverse image\index{Inverse image} $f^{-1}(X')$ of $X' \subset S'$ consists of all $x \in S$ such that $f(x) \in X'$. Observe that $f(f^{-1}(X')) \subset X'$, and $f^{-1}(f(X)) \supset X$.

The definition of a continuous function\index{Continuous function|(}\index{Function!continuous|(}\index{Continuous function}\index{Function!continuous} needs practically no modification: $f$ is continuous at $a$ if to every $\varepsilon > 0$ there exists $\delta > 0$ such that $d(x, a) < \delta$ implies $d'(f(x), f(a)) < \varepsilon$. We are mainly concerned with functions that are continuous at all points in the domain of definition.


% Original PDF physical page 80
\index{Compactness|)}\index{Into, definition|)}\index{Set!compact|)}\index{Weierstrass K@Weierstrass, K.|)}\sourcepage{80}
The following characterizations are immediate consequences of the definition:

A function is continuous if and only if the inverse image of every open set is open.

A function is continuous if and only if the inverse image of every closed set is closed.

If f is not defined on all of S, the words “open” and “closed,” when referring to the inverse image, should of course be interpreted relatively to the domain of f. It is very important to observe that these properties hold only for the inverse image, not for the direct image. For instance the mapping $f(x) = x^{2}/(1 + x^{2})$ of R into R has the image $f(\mathbf{R}) = \{y; 0 \leq y < 1\}$ which is neither open nor closed. In this example $f(\mathbf{R})$ fails to be closed because R is not compact. In fact, the following is true:

\begin{originalstatement}{Theorem}{8}\label{thm:3.8}
Under a continuous mapping\index{Mapping!continuous|(}\index{Mapping!continuous} the image of every compact set is compact, and consequently closed.
\end{originalstatement}
Suppose that $f$ is defined and continuous on the compact set $X$. Consider a covering of $f(X)$ by open sets $U$. The inverse images $f^{-1}(U)$ are open and form a covering of $X$. Because $X$ is compact we can select a finite subcovering: $X \subset f^{-1}(U_1) \cup \dots \cup f^{-1}(U_m)$. It follows that $f(X) \subset U_1 \cup \dots \cup U_m$, and we have proved that $f(X)$ is compact.

\begin{originalstatement}{Corollary}{}\label{cor:3.unnumbered-p80-b6}
A continuous real-valued function on a compact set has a maximum and a minimum.
\end{originalstatement}
The image is a closed bounded subset of R. The existence of a maximum and a minimum follows by Theorem~\hyperref[thm:3.2]{2}.

\begin{originalstatement}{Theorem}{9}\label{thm:3.9}
Under a continuous mapping the image of any connected set is connected.
\end{originalstatement}
We may assume that $f$ is defined and continuous on the whole space $S$, and that $f(S)$ is all of $S'$. Suppose that $S' = A \cup B$ where $A$ and $B$ are open and disjoint. Then $S = f^{-1}(A) \cup f^{-1}(B)$ is a representation of $S$ as a union of disjoint open sets. If $S$ is connected either $f^{-1}(A) = 0$ or $f^{-1}(B) = 0$, and hence $A = 0$ or $B = 0$. We conclude that $S'$ is connected.

A typical application is the assertion that a real-valued function which is continuous and never zero on a connected set is either always positive or always negative. In fact, the image is connected, and hence an interval. But an interval which contains positive and negative num%
% Original PDF physical page 81
\sourcepage{81}%
bers\index{Bijective, definition|(}\index{Continuous function!uniformly|(}\index{Injective, definition|(}\index{Inverse function|(}\index{One to one, definition|(}\index{Onto, definition|(}\index{Surjective, definition|(} also contains zero.

A mapping $f: S \to S'$ is said to be one to one if $f(x) = f(y)$ only for $x = y$; it is said to be onto if $f(S) = S'$.\hyperlink{fn:81:dagger}{\textsuperscript{\ensuremath{\dagger}}} A mapping with both these properties has an inverse $f^{-1}$, defined on $S'$; it satisfies $f^{-1}(f(x)) = x$ and $f(f^{-1}(x')) = x'$. In this situation, if $f$ and $f^{-1}$ are both continuous we say that $f$ is a topological mapping\index{Mapping!topological}\index{Topological mapping} or a homeomorphism. A property of a set which is shared by all topological images is called a topological property\index{Topological property}. For instance, we have proved that compactness and connectedness are topological properties (Theorems~\hyperref[thm:3.8]{8} and \hyperref[thm:3.9]{9}). In this connection it is perhaps useful to point out that the property of being an open subset is not topological. If $X \subset S$ and $Y \subset S'$ and if $X$ is homeomorphic to $Y$ there is no reason why $X$ and $Y$ should be simultaneously open. It happens to be true if $S = S' = \mathbf{R}^n$ (invariance of the region), but this is a deep theorem that we shall not need.

The notion of uniform continuity will be in constant use. Quite generally, a condition is said to hold uniformly with respect to a parameter if it can be expressed by inequalities which do not involve the parameter. Accordingly, a function f is said to be uniformly continuous on X if, to every $\varepsilon > 0$, there exists a $\delta > 0$ such that $d'(f(x_1), f(x_2)) < \varepsilon$ for all pairs $(x_1, x_2)$ with $d(x_1, x_2) < \delta$. The emphasis is on the fact that $\delta$ is not allowed to depend on $x_1$.

\begin{originalstatement}{Theorem}{10}\label{thm:3.10}
On a compact set every continuous function\index{Continuous function}\index{Function!continuous} is uniformly continuous.
\end{originalstatement}
The proof is typical of the way the Heine-Borel property can be used. Suppose that $f$ is continuous on a compact set $X$. For every $y \in X$ there is a ball $B(y, \rho)$ such that $d'(f(x), f(y)) < \epsilon/2$ for $x \in B(y, \rho)$; here $\rho$ may depend on $y$. Consider the covering of $X$ by the smaller balls $B(y, \rho/2)$. There exists a finite subcovering: $X \subset B(y_1, \rho_1/2) \cup \cdots \cup B(y_m, \rho_m/2)$. Let $\delta$ be the smallest of the numbers $\rho_1/2, \ldots, \rho_m/2$, and suppose that $d(x_1, x_2) < \delta$. There is a $y_k$ with $d(x_1, y_k) < \rho_k/2$, and we obtain $d(x_2, y_k) < \rho_k/2 + \delta \leq \delta_k$. Hence $d'(f(x_1), f(y_k)) < \epsilon/2$ and $d'(f(x_2), f(y_k)) < \epsilon/2$ so that $d'(f(x_1), f(x_2)) < \epsilon$ as desired.

On sets which are not compact some continuous functions are uniformly continuous and others are not. For instance, the function z is uniformly continuous on the whole complex plane, but the function $z^{2}$ is not.

\begingroup\renewcommand{\thefootnote}{}\footnotetext{\hypertarget{fn:81:dagger}{}\textsuperscript{\ensuremath{\dagger}} These linguistically clumsy terms can be replaced by \emph{injective} (for one to one) and \emph{surjective} (for onto). A mapping with both properties is called \emph{bijective}.}\endgroup

% Original PDF physical page 82
\index{Bijective, definition|)}\index{Continuous function!uniformly|)}\index{Injective, definition|)}\index{Inverse function|)}\index{One to one, definition|)}\index{Onto, definition|)}\index{Surjective, definition|)}\sourcepage{82}
\subsection*{EXERCISES}\index{Uniform continuity|(}
\originalnumber{1}{exercise:3.1.5.4.1}
1. Construct a topological mapping of the open disk $|z| < 1$ onto the whole plane.

\originalnumber{2}{exercise:3.1.5.4.2}
2. Prove that a subset of the real line which is topologically equivalent to an open interval is an open interval. (Consider the effect of removing a point.)

\originalnumber{3}{exercise:3.1.5.4.3}
3. Prove that every continuous one-to-one mapping of a compact space is topological. (Show that closed sets are mapped on closed sets.)

\originalnumber{4}{exercise:3.1.5.4.4}
4. Let $X$ and $Y$ be compact sets in a complete metric space. Prove that there exist $x \in X$, $y \in Y$ such that $d(x, y)$ is a minimum.

\originalnumber{5}{exercise:3.1.5.4.5}
5. Which of the following functions are uniformly continuous on the whole real line: $\sin x, x \sin x, x \sin (x^2), |x|^{\frac{1}{2}} \sin x$?

\setcounter{section}{1}\setcounter{subsection}{5}\subsection{Topological Spaces}\label{sec:3.1.6}
It is not necessary, and not always convenient, to express nearness in terms of distance. The observant reader will have noticed that most results in the preceding sections were formulated in terms of open sets. True enough, we used distances to define open sets, but there is really no strong reason to do this. If we decide to consider the open sets as the primary objects we must postulate axioms that they have to satisfy. The following axioms lead to the commonly accepted definition of a topological space:

\begin{originalstatement}{Definition}{8}\label{def:3.8}
A topological space is a set T together with a collection of its subsets, called open sets. The following conditions have to be fulfilled:
(i) The empty set $\varnothing$ and the whole space $T$ are open sets.

(ii) The intersection of any two open sets is an open set.

(iii) The union of an arbitrary collection of open sets is an open set.
\end{originalstatement}

We recognize at once that this terminology is consistent with our earlier definition of an open subset of a metric space. Indeed, properties (ii) and (iii) were strongly emphasized, and (i) is trivial.

Closed sets are the complements of open sets, and it is immediately clear how to define interior, closure, boundary, and so on. Neighborhoods could be avoided, but they are rather convenient: N is a neighborhood of x if there exists an open set U such that $x \in U$ and $U \subset N$.

Connectedness was defined purely by means of open sets. Hence the definition carries over to topological spaces, and the theorems remain true. The Heine-Borel property is also one that deals only with open sets. Therefore it makes perfect sense to speak of a compact topological space. However, Theorem~\hyperref[thm:3.6]{6} becomes meaningless, and Theorem~\hyperref[thm:3.7]{7} becomes false.

As a matter of fact, the first serious difficulty we encounter is with
% Original PDF physical page 83
\index{Continuous function|)}\index{Function!continuous|)}\index{Uniform continuity|)}\sourcepage{83}
convergent\index{Arc|(}\index{Conformal mapping|(} sequences. The definition is clear: we say that $x_{n} \rightarrow x$ if every neighborhood of x contains all but a finite number of the $x_{n}$. But if $x_{n} \rightarrow x$ and $x_{n} \rightarrow y$ we are not able to prove that x = y. This awkward situation is remedied by introducing a new axiom which characterizes the topological space\index{Space!topological|(}\index{Space!topological} as a Hausdorff space\index{Space!Hausdorff}:

\begin{originalstatement}{Definition}{9}\label{def:3.9}
A topological space is called a Hausdorff space if any two distinct points are contained in disjoint open sets.
\end{originalstatement}
In other words, if $x \neq y$ we require the existence of open sets U, V such that $x \in U$, $y \in V$ and $U \cap V = \varnothing$. In the presence of this condition it is obvious that the limit of a convergent sequence is unique. We shall never in this book have occasion to consider a space that is not a Hausdorff space.

This is not the place to give examples of topologies that cannot be derived from a distance function. Such examples would necessarily be very complicated and would not further the purposes of this book. The point is that it may be unnatural to introduce a distance in situations when one is not really needed. The reason for including this section has been to alert the reader that distances are dispensable.

\setcounter{section}{1}\section{Conformality}\label{sec:3.2}
We now return to our original setting where all functions and variables are restricted to real or complex numbers. The role of metric spaces will seem disproportionately small: all we actually need are some simple applications of connectedness and compactness.

The whole section is mainly descriptive. It centers on the geometric consequences of the existence of a derivative.

\setcounter{section}{2}\setcounter{subsection}{0}\subsection{Arcs and Closed Curves}\label{sec:3.2.1}
The equation of an arc $\gamma$ in the plane is most conveniently given in parametric form $x = x(t)$, $y = y(t)$ where t runs through an interval $\alpha \leq t \leq \beta$ and $x(t)$, $y(t)$ are continuous functions. We can also use the complex notation $z = z(t) = x(t) + iy(t)$ which has several advantages. It is also customary to identify the arc $\gamma$ with the continuous mapping\index{Mapping!continuous} of $[\alpha, \beta]$. When following this custom it is preferable to denote the mapping by $z = \gamma(t)$.

Considered as a point set an arc is the image of a closed finite interval under a continuous mapping. As such it is compact and connected. However, an arc is not merely a set of points, but very essentially also a succession of points, ordered by increasing values of the parameter. If a nondecreasing function $t = \varphi(\tau)$ maps an interval $\alpha' \leq \tau \leq \beta'$ onto $\alpha \leq t \leq \beta$, then $z = z(\varphi(\tau))$ defines the same succession of points as $z = z(t)$.


% Original PDF physical page 84
\index{Mapping!continuous|)}\sourcepage{84}
We\index{Arc!differentiable|(}\index{Arc!simple|(}\index{Differentiable arc|(}\index{Piecewise differentiable arc|(} say that the first equation arises from the second by a change of parameter\index{Change of parameter}\index{Parameter!change of}\index{Parameter}. The change is reversible if and only if $\varphi(\tau)$ is strictly increasing. For instance, the equation $z = t^{2} + it^{4}, 0 \leq t \leq 1$ arises by a reversible change of parameter\index{Change of parameter!reversible} from the equation $z = t + it^{2}, 0 \leq t \leq 1$. A change of the parametric interval $(\alpha, \beta)$ can always be brought about by a linear change of parameter, which is one of the form $t = a\tau + b, a > 0$.

Logically, the simplest course is to consider two arcs as different as soon as they are given by different equations, regardless of whether one equation may arise from the other by a change of parameter. In following this course, as we will, it is important to show that certain properties of arcs are invariant under a change of parameter. For instance, the initial and terminal point of an arc remain the same after a change of parameter.

If the derivative $z'(t) = x'(t) + iy'(t)$ exists and is $\neq 0$, the arc $\gamma$ has a tangent whose direction is determined by $\arg z'(t)$. We shall say that the arc is differentiable if $z'(t)$ exists and is continuous (the term continuously differentiable is too unwieldy); if, in addition, $z'(t) \neq 0$ the arc is said to be regular. An arc is piecewise differentiable or piecewise regular if the same conditions hold except for a finite number of values t; at these points $z(t)$ shall still be continuous with left and right derivatives which are equal to the left and right limits of $z'(t)$ and, in the case of a piecewise regular arc\index{Arc!regular}\index{Regular arc}, $\neq 0$.

The differentiable or regular character of an arc is invariant under the change of parameter $t = \varphi(\tau)$ provided that $\varphi'(\tau)$ is continuous and, for regularity, $\neq0$. When this is the case, we speak of a differentiable or regular change of parameter.

An arc is simple, or a Jordan arc\index{Arc!Jordan}\index{Jordan arc}, if $z(t_{1}) = z(t_{2})$ only for $t_{1} = t_{2}$. An arc is a closed curve\index{Closed curve}\index{Curve} if the end points coincide: $z(\alpha) = z(\beta)$. For closed curves a shift of the parameter is defined as follows: If the original equation is $z = z(t)$, $\alpha \leq t \leq \beta$, we choose a point $t_{0}$ from the interval $(\alpha, \beta)$ and define a new closed curve whose equation is $z = z(t)$ for $t_{0} \leq t \leq \beta$ and $z = z(t - \beta + \alpha)$ for $\beta \leq t \leq t_{0} + \beta - \alpha$. The purpose of the shift is to get rid of the distinguished position of the initial point. The correct definitions of a differentiable or regular closed curve and of a simple closed curve (or Jordan curve\index{Curve!Jordan}\index{Jordan curve}) are obvious.

The opposite arc\index{Arc!opposite} of $z = z(t)$, $\alpha \leq t \leq \beta$, is the arc $z = z(-t)$, $-\beta \leq t \leq -\alpha$. Opposite arcs are sometimes denoted by $\gamma$ and $-\gamma$, sometimes by $\gamma$ and $\gamma^{-1}$, depending on the connection. A constant function $z(t)$ defines a point curve\index{Curve!point}\index{Point curve}.

A circle $C$, originally defined as a locus $|z - a| = r$, can be considered as a closed curve with the equation $z = a + re^{it}$, $0 \leq t \leq 2\pi$. We will use this standard parametrization whenever a circle is introduced.

This convention saves us from writing down the equation each time it is
% Original PDF physical page 85
\index{Arc!differentiable|)}\index{Arc!simple|)}\index{Differentiable arc|)}\index{Piecewise differentiable arc|)}\index{Space!topological|)}\sourcepage{85}
needed;\index{Tangent|(} also, and this is its most important purpose, it serves as a definite rule to distinguish between C and -C.

\setcounter{section}{2}\setcounter{subsection}{1}\subsection{Analytic Functions in Regions}\label{sec:3.2.2}
When we consider the derivative
\begin{equation*}
f ^ {\prime} (z) = \lim _ {h \rightarrow 0} \frac {f (z + h) - f (z)}{h}
\end{equation*}
of a complex-valued function, defined on a set A in the complex plane, it is of course understood that $z \in A$ and that the limit is with respect to values h such that $z + h \in A$. The existence of the derivative will therefore have a different meaning depending on whether z is an interior point or a boundary point of A. The way to avoid this is to insist that all analytic functions be defined on open sets.

We give a formal statement of the definition:

\begin{originalstatement}{Definition}{10}\label{def:3.10}
A complex-valued function $f(z)$, defined on an open set $\Omega$, is said to be analytic in $\Omega$ if it has a derivative at each point of $\Omega$.
\end{originalstatement}
Sometimes one says more explicitly that $f(z)$ is complex analytic. A commonly used synonym is holomorphic.

It is important to stress that the open set $\Omega$ is part of the definition. As a rule one should avoid speaking of an analytic function $f(z)$ without referring to a specific open set $\Omega$ on which it is defined, but the rule can be broken if it is clear from the context what the set is. Observe that f must first of all be a function, and hence single-valued. If $\Omega'$ is an open subset of $\Omega$, and if $f(z)$ is analytic in $\Omega$, then the restriction of f to $\Omega'$ is analytic in $\Omega'$; it is customary to denote the restriction by the same letter f. In particular, since the components of an open set are open, it is no loss of generality to consider only the case where $\Omega$ is connected, that is to say a region.

For greater flexibility of the language it is desirable to introduce the following complement to Definition~\hyperref[def:3.10]{10}:

\begin{originalstatement}{Definition}{11}\label{def:3.11}
A function $f(z)$ is analytic on an arbitrary point set $A$ if it is the restriction to $A$ of a function which is analytic in some open set containing $A$.
\end{originalstatement}
The last definition is merely an agreement to use a convenient terminology. This is a case in which the set $\Omega$ need not be explicitly mentioned, for the specific choice of $\Omega$ is usually immaterial as long as it contains A. Another instance in which the mention of $\Omega$ can be suppressed is the phrase: “Let $f(z)$ be analytic at $z_{0}$.” It means that a function $f(z)$ is defined and has a derivative in some unspecified open neighborhood of $z_{0}$.


% Original PDF physical page 86
\index{Arc|)}\index{Tangent|)}\sourcepage{86}
Although our definition requires all analytic functions to be single-valued, it is possible to consider such multiple-valued functions as $\sqrt{z}$, $\log z$, or arc cos z, provided that they are restricted to a definite region in which it is possible to select a single-valued and analytic branch of the function.

For instance, we may choose for $\Omega$ the complement of the negative real axis $z \leq 0$; this set is indeed open and connected. In $\Omega$ one and only one of the values of $\sqrt{z}$ has a positive real part. With this choice $w = \sqrt{z}$ becomes a single-valued function in $\Omega$; let us prove that it is continuous. Choose two points $z_1, z_2 \in \Omega$ and denote the corresponding values of $w$ by $w_1 = u_1 + iv_1$, $w_2 = u_2 + iv_2$ with $u_1, u_2 > 0$. Then
\begin{equation*}
\left| z _ {1} - z _ {2} \right| = \left| w _ {1} ^ {2} - w _ {2} ^ {2} \right| = \left| w _ {1} - w _ {2} \right| \cdot \left| w _ {1} + w _ {2} \right|
\end{equation*}
and $|w_{1} + w_{2}| \geq u_{1} + u_{2} > u_{1}$. Hence
\begin{equation*}
\left| w _ {1} - w _ {2} \right| <   \frac {\left| z _ {1} - z _ {2} \right|}{u _ {1}}
\end{equation*}
and it follows that $w = \sqrt{z}$ is continuous at $z_{1}$. Once the continuity is established the analyticity follows by derivation of the inverse function $z = w^{2}$. Indeed, with the notations used in calculus $\Delta z \to 0$ implies $\Delta w \to 0$. Therefore,
\begin{equation*}
\lim _ {\Delta z \to 0} \frac {\Delta w}{\Delta z} = \lim _ {\Delta w \to 0} \frac {\Delta w}{\Delta z}
\end{equation*}
and we obtain
\begin{equation*}
{\frac {d w}{d z}} = {\frac {1}{\frac {d z}{d w}}} = {\frac {1}{2 w}} = {\frac {1}{2 \sqrt {z}}}
\end{equation*}
with the same branch of $\sqrt{z}$.

In the case of $\log z$ we can use the same region $\Omega$, obtained by excluding the negative real axis, and define the principal branch of the logarithm by the condition $|\operatorname{Im} \log z| < \pi$. Again, the continuity must be proved, but this time we have no algebraic identity at our disposal, and we are forced to use a more general reasoning. Denote the principal branch by $w = u + iv = \log z$. For a given point $w_1 = u_1 + iv_1$, $|v_1| < \pi$, and a given $\varepsilon > 0$, consider the set $A$ in the $w$-plane which is defined by the inequalities $|w - w_1| \geq \varepsilon$, $|v| \leq \pi$, $|u - u_1| \leq \log 2$. This set is closed and bounded, and for sufficiently small $\varepsilon$ it is not empty. The continuous function $|e^w - e^{w_1}|$ has consequently a minimum $\rho$ on $A$ (Theorem~\hyperref[thm:3.8]{8}, Corollary). This minimum is positive, for $A$ does not contain any point $w_1 + n \cdot 2\pi i$. Choose $\delta = \min (\rho, \frac{1}{2} e^{u_1})$, and assume that
\begin{equation*}
\left| z _ {1} - z _ {2} \right| = \left| e ^ {w _ {1}} - e ^ {w _ {2}} \right| <   \delta .
\end{equation*}
% Original PDF physical page 87
\sourcepage{87}
Then\index{Principal branch|(} $w_{2}$ cannot lie in $A$, for this would make $|e^{w_1} - e^{w_2}| \geq \rho \geq \delta$. Neither is it possible that $u_{2} < u_{1} - \log 2$ or $u_{2} > u_{1} + \log 2$; in the former case we would obtain $|e^{w_1} - e^{w_2}| \geq e^{u_1} - e^{u_2} > \frac{1}{2} e^{u_1} \geq \delta$, and in the latter case $|e^{w_1} - e^{w_2}| \geq e^{u_2} - e^{u_1} > e^{u_1} > \delta$. Hence $w_{2}$ must lie in the disk $|w - w_1| < \varepsilon$, and we have proved that $w$ is a continuous function of $z$. From the continuity we conclude as above that the derivative exists and equals $1/z$.

The infinitely many values of arc cos z are the same as the values of $i \log (z + \sqrt{z^{2} - 1})$. In this case we restrict z to the complement $\Omega'$ of the half lines $x \leq -1$, y = 0 and $x \geq 1$, y = 0. Since $1 - z^{2}$ is never real and $\leq 0$ in $\Omega'$, we can define $\sqrt{1 - z^{2}}$ as in the first example and then set $\sqrt{z^{2} - 1} = i\sqrt{1 - z^{2}}$. Moreover, $z + \sqrt{z^{2} - 1}$ is never real in $\Omega'$, for $z + \sqrt{z^{2} - 1}$ and $z - \sqrt{z^{2} - 1}$ are reciprocals and hence real only if z and $\sqrt{z^{2} - 1}$ are both real; this happens only when z lies on the excluded parts of the real axis. Because $\Omega'$ is connected, it follows that all values of $z + \sqrt{z^{2} - 1}$ in $\Omega'$ are on the same side of the real axis, and since i is such a value they are all in the upper half plane. We can therefore define an analytic branch of $\log (z + \sqrt{z^{2} - 1})$ whose imaginary part lies between 0 and $\pi$. In this way we obtain a single-valued analytic function
\begin{equation*}
\operatorname{arc} \cos z = i \log (z + \sqrt {z ^ {2} - 1})
\end{equation*}
in $\Omega^{\prime}$ whose derivative is
\begin{equation*}
D \operatorname{arc} \cos z = i \frac {1}{z + \sqrt {z ^ {2} - 1}} \left(1 + \frac {z}{\sqrt {z ^ {2} - 1}}\right) = \frac {1}{\sqrt {1 - z ^ {2}}}
\end{equation*}
where $\sqrt{1 - z^2}$ has a positive real part.

There is nothing unique about the way in which the region and the single-valued branches have been chosen in these examples. Therefore, each time we consider a function such as $\log z$ the choice of the branch has to be specified. It is a fundamental fact that it is impossible to define a single-valued and analytic branch of $\log z$ in certain regions. This will be proved in the chapter on integration.

All the results of Chap. II, Sec.~\hyperref[sec:3.1.2]{1.2} remain valid for functions which are analytic on an open set. In particular, the real and imaginary parts of an analytic function in $\Omega$ satisfy the Cauchy-Riemann equations
\begin{equation*}
\frac {\partial u}{\partial x} = \frac {\partial v}{\partial y}, \quad \frac {\partial u}{\partial y} = - \frac {\partial v}{\partial x}.
\end{equation*}
Conversely, if u and v satisfy these equations in $\Omega$, and if the partial derivatives are continuous, then $u + iv$ is an analytic function in $\Omega$.

An analytic function in $\Omega$ degenerates if it reduces to a constant. In
% Original PDF physical page 88
\index{Principal branch|)}\sourcepage{88}
the following theorem we shall list some simple conditions which have this consequence:

\begin{originalstatement}{Theorem}{11}\label{thm:3.11}
An analytic function in a region $\Omega$ whose derivative vanishes identically must reduce to a constant. The same is true if either the real part, the imaginary part, the modulus, or the argument is constant.
\end{originalstatement}
The vanishing of the derivative implies that $\partial u/\partial x$, $\partial u/\partial y$, $\partial v/\partial x$, $\partial v/\partial y$ are all zero. It follows that u and v are constant on any line segment in $\Omega$ which is parallel to one of the coordinate axes. In Sec.~\hyperref[sec:3.1.3]{1.3} we remarked, in connection with Theorem~\hyperref[thm:3.3]{3}, that any two points in a region can be joined within the region by a polygon whose sides are parallel to the axes. We conclude that $u + iv$ is constant.

If $u$ or $v$ is constant,
\begin{equation*}
f ^ {\prime} (z) = \frac {\partial u}{\partial x} - i \frac {\partial u}{\partial y} = \frac {\partial v}{\partial y} + i \frac {\partial v}{\partial x} = 0,
\end{equation*}
and hence $f(z)$ must be constant. If $u^2 + v^2$ is constant, we obtain
\begin{equation*}
u \frac {\partial u}{\partial x} + v \frac {\partial v}{\partial x} = 0
\end{equation*}
and
\begin{equation*}
u \frac {\partial u}{\partial y} + v \frac {\partial v}{\partial y} = - u \frac {\partial v}{\partial x} + v \frac {\partial u}{\partial x} = 0.
\end{equation*}
These equations permit the conclusion $\partial u/\partial x = \partial v/\partial x = 0$ unless the determinant $u^{2} + v^{2}$ vanishes. But if $u^{2} + v^{2} = 0$ at a single point it is constantly zero and $f(z)$ vanishes identically. Hence $f(z)$ is in any case a constant.

Finally, if $\arg f(z)$ is constant, we can set u = kv with constant k (unless v is identically zero). But u - kv is the real part of $(1 + ik)f$, and we conclude again that f must reduce to a constant.

Note that for this theorem it is essential that $\Omega$ is a region. If not, we can only assert that $f(z)$ is constant on each component of $\Omega$.

\subsection*{EXERCISES}
\originalnumber{1}{exercise:3.2.2.5.1}
1. Give a precise definition of a single-valued branch of $\sqrt{1 + z} + \sqrt{1 - z}$ in a suitable region, and prove that it is analytic.

\originalnumber{2}{exercise:3.2.2.5.2}
2. Same problem for log log z.

\originalnumber{3}{exercise:3.2.2.5.3}
3. Suppose that $f(z)$ is analytic and satisfies the condition $|f(z)^{2}-1|<1$ in a region $\Omega$. Show that either $\operatorname{Re} f(z)>0$ or $\operatorname{Re} f(z)<0$ throughout $\Omega$.


% Original PDF physical page 89
\sourcepage{89}
\setcounter{section}{2}\setcounter{subsection}{2}\subsection{Conformal Mapping}\index{Mapping!conformal|(}\index{Conformal mapping}\index{Mapping!conformal}\label{sec:3.2.3}
Suppose that an arc $\gamma$ with the equation $z = z(t)$, $\alpha \leq t \leq \beta$, is contained in a region $\Omega$, and let $f(z)$ be defined and continuous in $\Omega$. Then the equation $w = w(t) = f(z(t))$ defines an arc $\gamma'$ in the w-plane which may be called the image\index{Image} of $\gamma$.

Consider the case of an $f(z)$ which is analytic in $\Omega$. If $z'(t)$ exists, we find that $w'(t)$ also exists and is determined by
\begin{equation*}
w ^ {\prime} (t) = f ^ {\prime} (z (t)) z ^ {\prime} (t).\tag{1}\label{eq:3.1}
\end{equation*}
We will investigate the meaning of this equation at a point $z_0 = z(t_0)$ with $z'(t_0) \neq 0$ and $f'(z_0) \neq 0$.

The first conclusion is that $w'(t_0) \neq 0$. Hence $\gamma'$ has a tangent at $w_0 = f(z_0)$, and its direction is determined by
\begin{equation*}
\arg w ^ {\prime} (t _ {0}) = \arg f ^ {\prime} (z _ {0}) + \arg z ^ {\prime} (t _ {0}).\tag{2}\label{eq:3.2}
\end{equation*}
This relation asserts that the angle between the directed tangents to $\gamma$ at $z_{0}$ and to $\gamma'$ at $w_{0}$ is equal to $\arg f'(z_{0})$. It is hence independent of the curve $\gamma$. For this reason curves through $z_{0}$ which are tangent to each other are mapped onto curves with a common tangent at $w_{0}$. Moreover, two curves which form an angle at $z_{0}$ are mapped upon curves forming the same angle, in sense as well as in size. In view of this property the mapping by $w = f(z)$ is said to be conformal at all points with $f'(z) \neq 0$.

A related property of the mapping is derived by consideration of the modulus $|f'(z_{0})|$. We have
\begin{equation*}
\lim _ {z \rightarrow z _ {0}} \frac {| f (z) - f (z _ {0}) |}{| z - z _ {0} |} = | f ^ {\prime} (z _ {0}) |,
\end{equation*}
and this means that any small line segment with one end point at $z_{0}$ is, in the limit, contracted or expanded in the ratio $|f'(z_{0})|$. In other words, the linear change of scale at $z_{0}$, effected by the transformation $w = f(z)$, is independent of the direction. In general this change of scale will vary from point to point.

Conversely, it is clear that both kinds of conformality together imply the existence of $f'(z_{0})$. It is less obvious that each kind will separately imply the same result, at least under additional regularity assumptions.

To be more precise, let us assume that the partial derivatives $\partial f/\partial x$ and $\partial f/\partial y$ are continuous. Under this condition the derivative of $w(t) = f(z(t))$ can be expressed in the form
\begin{equation*}
w ^ {\prime} (t _ {0}) = \frac {\partial f}{\partial x} x ^ {\prime} (t _ {0}) + \frac {\partial f}{\partial y} y ^ {\prime} (t _ {0})
\end{equation*}
where the partial derivatives are taken at $z_0$. In terms of $z'(t_0)$ this can


% Original PDF physical page 90
\sourcepage{90}
be\index{Indirectly conformal mapping|(} rewritten as
\begin{equation*}
w ^ {\prime} (t _ {0}) = \frac {1}{2} \left(\frac {\partial f}{\partial x} - i \frac {\partial f}{\partial y}\right) z ^ {\prime} (t _ {0}) + \frac {1}{2} \left(\frac {\partial f}{\partial x} + i \frac {\partial f}{\partial y}\right) \overline {{z ^ {\prime} (t _ {0})}}.
\end{equation*}
If angles are preserved, $\arg [w'(t_0) / z'(t_0)]$ must be independent of $\arg z'(t_0)$. The expression
\begin{equation*}
\frac {1}{2} \left(\frac {\partial f}{\partial x} - i \frac {\partial f}{\partial y}\right) + \frac {1}{2} \left(\frac {\partial f}{\partial x} + i \frac {\partial f}{\partial y}\right) \frac {\overline {{z ^ {\prime} (t _ {0})}}}{z ^ {\prime} (t _ {0})}\tag{3}\label{eq:3.3}
\end{equation*}
must therefore have a constant argument. As $\arg z'(t_{0})$ is allowed to vary, the point represented by \hyperref[eq:3.3]{(3)} describes a circle having the radius $\frac{1}{2}|(\partial f/\partial x)+i(\partial f/\partial y)|$. The argument cannot be constant on this circle unless its radius vanishes, and hence we must have
\begin{equation*}
\frac {\partial f}{\partial x} = - i \frac {\partial f}{\partial y}\tag{4}\label{eq:3.4}
\end{equation*}
which is the complex form of the Cauchy-Riemann equations.

Quite similarly, the condition that the change of scale shall be the same in all directions implies that the expression \hyperref[eq:3.3]{(3)} has a constant modulus. On a circle the modulus is constant only if the radius vanishes or if the center lies at the origin. In the first case we obtain \hyperref[eq:3.4]{(4)}, and in the second case
\begin{equation*}
\frac {\partial f}{\partial x} = i \frac {\partial f}{\partial y}.
\end{equation*}
The last equation expresses the fact that $\overline{f(z)}$ is analytic. A mapping by the conjugate of an analytic function with a nonvanishing derivative is said to be indirectly conformal. It evidently preserves the size but reverses the sense of angles.

If the mapping of $\Omega$ by $w = f(z)$ is topological, then the inverse function $z = f^{-1}(w)$ is also analytic. This follows easily if $f'(z) \neq 0$, for then the derivative of the inverse function must be equal to $1/f'(z)$ at the point $z = f^{-1}(w)$. We shall prove later that $f'(z)$ can never vanish in the case of a topological mapping by an analytic function.

The knowledge that $f'(z_0) \neq 0$ is sufficient to conclude that the mapping is topological if it is restricted to a sufficiently small neighborhood of $z_0$. This follows by the theorem on implicit functions known from the calculus, for the Jacobian\index{Jacobian} of the functions $u = u(x,y)$, $v = v(x,y)$ at the point $z_0$ is $|f'(z_0)|^2$ and hence $\neq 0$. Later we shall present a simpler proof of this important theorem.

But even if $f'(z) \neq 0$ throughout the region $\Omega$, we cannot assert that the mapping of the whole region is necessarily topological. To illustrate
% Original PDF physical page 91
\index{Indirectly conformal mapping|)}\sourcepage{91}
what may happen we refer to Fig.~\hyperref[fig:3-1]{3-1}. Here the mappings of the sub-

\bookfigure{91}{1}{FIG. 3-1. Doubly covered region.}
regions $\Omega_{1}$ and $\Omega_{2}$ are one to one, but the images overlap. It is helpful to think of the image of the whole region as a transparent film which partly covers itself. This is the simple and fruitful idea used by Riemann when he introduced the generalized regions now known as Riemann surfaces.

\setcounter{section}{2}\setcounter{subsection}{3}\subsection{Length and Area}\label{sec:3.2.4}
We have found that under a conformal mapping\index{Conformal mapping}\index{Mapping!conformal} $f(z)$ the length\index{Length} of an infinitesimal line segment at the point z is multiplied by the factor $|f'(z)|$. Because the distortion is the same in all directions, infinitesimal areas will clearly be multiplied by $|f'(z)|^{2}$.

Let us put this on a rigorous basis. We know from calculus that the length of a differentiable arc $\gamma$ with the equation $z = z(t) = x(t) + iy(t)$, $a \leq t \leq b$, is given by
\begin{equation*}
L (\gamma) = \int_ {a} ^ {b} \sqrt {x ^ {\prime} (t) ^ {2} + y ^ {\prime} (t) ^ {2}} d t = \int_ {a} ^ {b} | z ^ {\prime} (t) | d t.
\end{equation*}
The image curve $\gamma'$ is determined by $w = w(t) = f(z(t))$ with the derivative $w'(t) = f'(z(t))z'(t)$. Its length is thus
\begin{equation*}
L (\gamma^ {\prime}) = \int_ {a} ^ {b} | f ^ {\prime} (z (t)) | | z ^ {\prime} (t) | d t.
\end{equation*}
It is customary to use the shorter notations
\begin{equation*}
L (\gamma) = \int_ {\gamma} | d z |, \quad L (\gamma^ {\prime}) = \int_ {\gamma} | f ^ {\prime} (z) | | d z |.\tag{5}\label{eq:3.5}
\end{equation*}
Observe that in complex notation the calculus symbol ds for integration with respect to arc length\index{Arc length} is replaced by $|dz|$.

Now let E be a point set in the plane whose area\index{Area|(}\index{Area}
\begin{equation*}
A (E) = \iint_ {E} d x d y
\end{equation*}
% Original PDF physical page 92
\index{Mapping!conformal|)}\sourcepage{92}
can be evaluated as a double Riemann integral. If $f(z) = u(x,y) + iv(x,y)$ is a bijective differentiable mapping, then by the rule for changing integration variables the area\index{Area} of the image $E' = f(E)$ is given by
\begin{equation*}
A (E ^ {\prime}) = \iint_ {E} | u _ {x} v _ {y} - u _ {y} v _ {x} | d x d y.
\end{equation*}
But if $f(z)$ is a conformal mapping\index{Conformal mapping} of an open set containing E, then $u_{x}v_{y} - u_{y}v_{x} = |f'(z)|^{2}$ by virtue of the Cauchy-Riemann equations, and we obtain
\begin{equation*}
A (E ^ {\prime}) = \iint_ {E} | f ^ {\prime} (z) | ^ {2} d x d y.\tag{6}\label{eq:3.6}
\end{equation*}
The formulas \hyperref[eq:3.5]{(5)} and \hyperref[eq:3.6]{(6)} have important applications in the part of complex analysis that is frequently referred to as geometric function theory.

\setcounter{section}{2}\section{Linear Transformations}\label{sec:3.3}
Of all analytic functions the first-order rational functions have the simplest mapping properties, for they define mappings of the extended plane onto itself which are at the same time conformal and topological. The linear transformations\index{Linear transformation|(}\index{Linear transformation} have also very remarkable geometric properties, and for that reason their importance goes far beyond serving as simple examples of conformal mappings. The reader will do well to pay particular attention to this geometric aspect, for it will equip him with simple but very valuable techniques.

\setcounter{section}{3}\setcounter{subsection}{0}\subsection{The Linear Group}\index{Linear group|(}\index{Linear group}\label{sec:3.3.1}
We have already remarked in Chap.~\hyperref[chap:2]{2}, Sec.~\hyperref[sec:2.1.4]{1.4} that a linear fractional transformation
\begin{equation*}
w = S (z) = \frac {a z + b}{c z + d}\tag{7}\label{eq:3.7}
\end{equation*}
with $ad - bc \neq 0$ has an inverse
\begin{equation*}
z = S ^ {- 1} (w) = \frac {d w - b}{- c w + a}.
\end{equation*}
The special values $S(\infty) = a/c$ and $S(-d/c) = \infty$ can be introduced either by convention or as limits for $z \to \infty$ and $z \to -d/c$. With the latter interpretation it becomes obvious that S is a topological mapping of the extended plane onto itself, the topology being defined by distances on the Riemann sphere.

For linear transformations we shall usually replace the notation $S(z)$


% Original PDF physical page 93
\index{Area|)}\sourcepage{93}
by Sz. The representation \hyperref[eq:3.7]{(7)} is said to be normalized if $ad - bc = 1$. It is clear that every linear transformation\index{Linear transformation} has two normalized representations, obtained from each other by changing the signs of the coefficients.

A convenient way to express a linear transformation is by use of homogeneous coordinates. If we write $z = z_{1}/z_{2}$, $w = w_{1}/w_{2}$ we find that w = Sz if
\begin{equation*}
\begin{array}{l} w _ {1} = a z _ {1} + b z _ {2} \\ w _ {2} = c z _ {1} + d z _ {2} \end{array}\tag{8}\label{eq:3.8}
\end{equation*}
or, in matrix notation,
\begin{equation*}
\binom {w _ {1}} {w _ {2}} = \left( \begin{array}{c c} a & b \\ c & d \end{array} \right) \binom {z _ {1}} {z _ {2}}.
\end{equation*}
The main advantage of this notation is that it leads to a simple determination of a composite transformation $w = S_{1}S_{2}z$. If we use subscripts to distinguish between the matrices that correspond to $S_{1}, S_{2}$ it is immediate that $S_{1}S_{2}$ belongs to the matrix product
\begin{equation*}
\left( \begin{array}{c c} a _ {1} & b _ {1} \\ c _ {1} & d _ {1} \end{array} \right) \left( \begin{array}{c c} a _ {2} & b _ {2} \\ c _ {2} & d _ {2} \end{array} \right) = \left( \begin{array}{c c} a _ {1} a _ {2} + b _ {1} c _ {2} & a _ {1} b _ {2} + b _ {1} d _ {2} \\ c _ {1} a _ {2} + d _ {1} c _ {2} & c _ {1} b _ {2} + d _ {1} d _ {2} \end{array} \right).
\end{equation*}
All linear transformations form a group. Indeed, the associative law $(S_{1}S_{2})S_{3}=S_{1}(S_{2}S_{3})$ holds for arbitrary transformations, the identity w=z is a linear transformation, and the inverse of a linear transformation is linear. The ratios $z_{1}:z_{2}\neq0:0$ are the points of the complex projective line, and \hyperref[eq:3.8]{(8)} identifies the group of linear transformations with the one-dimensional projective group over the complex numbers, usually denoted by $P(1,\mathbf{C})$. If we use only normalized representations, we can also identify it with the group of two-by-two matrices with determinant 1 (denoted $SL(2,\mathbf{C})$ ), except that there are two opposite matrices corresponding to the same linear transformation.

We shall make no further use of the matrix notation, except for remarking that the simplest linear transformations belong to matrices of the form
\begin{equation*}
\left( \begin{array}{c c} 1 & \alpha \\ 0 & 1 \end{array} \right), \left( \begin{array}{c c} k & 0 \\ 0 & 1 \end{array} \right), \left( \begin{array}{c c} 0 & 1 \\ 1 & 0 \end{array} \right).
\end{equation*}
The first of these, $w = z + \alpha$, is called a parallel translation. The second, w = kz, is a rotation if $|k| = 1$ and a homothetic transformation\index{Homothetic transformation} if k > 0. For arbitrary complex $k \neq 0$ we can set $k = |k| \cdot k / |k|$, and hence w = kz can be represented as the result of a homothetic transformation followed by a rotation. The third transformation, w = 1/z, is called an inversion\index{Inversion}.

If $c \neq 0$ we can write
\begin{equation*}
\frac {a z + b}{c z + d} = \frac {b c - a d}{c ^ {2} (z + d / c)} + \frac {a}{c},
\end{equation*}
% Original PDF physical page 94
\sourcepage{94}
and this decomposition shows that the most general linear transformation\index{Linear transformation} is composed by a translation, an inversion, a rotation\index{Rotation}, and a homothetic transformation followed by another translation. If c = 0, the inversion falls out and the last translation is not needed.

\subsection*{EXERCISES}
\originalnumber{1}{exercise:3.3.1.6.1}
1. Prove that the reflection $z \rightarrow \bar{z}$ is not a linear transformation.

\originalnumber{2}{exercise:3.3.1.6.2}
2. If
\begin{equation*}
T _ {1} z = \frac {z + 2}{z + 3}, \quad T _ {2} z = \frac {z}{z + 1},
\end{equation*}
find $T_{1}T_{2}z$, $T_{2}T_{1}z$ and $T_{1}^{-1}T_{2}z$.

\originalnumber{3}{exercise:3.3.1.6.3}
3. Prove that the most general transformation which leaves the origin fixed and preserves all distances is either a rotation or a rotation followed by reflexion in the real axis.

\originalnumber{4}{exercise:3.3.1.6.4}
4. Show that any linear transformation which transforms the real axis into itself can be written with real coefficients.

\setcounter{section}{3}\setcounter{subsection}{1}\subsection{The Cross Ratio}\label{sec:3.3.2}
Given three distinct points $z_{2}$, $z_{3}$, $z_{4}$ in the extended plane, there exists a linear transformation S which carries them into 1, 0, $\infty$ in this order. If none of the points is $\infty$, S will be given by
\begin{equation*}
S z = \frac {z - z _ {3}}{z - z _ {4}}: \frac {z _ {2} - z _ {3}}{z _ {2} - z _ {4}}.\tag{9}\label{eq:3.9}
\end{equation*}
If $z_{2}, z_{3}$ or $z_{4} = \infty$ the transformation reduces to
\begin{equation*}
\frac {z - z _ {3}}{z - z _ {4}}, \quad \frac {z _ {2} - z _ {4}}{z - z _ {4}}, \quad \frac {z - z _ {3}}{z _ {2} - z _ {3}}
\end{equation*}
respectively.

If T were another linear transformation with the same property, then $ST^{-1}$ would leave 1, 0, $\infty$ invariant. Direct calculation shows that this is true only for the identity transformation, and we would have S = T. We conclude that S is uniquely determined.

\begin{originalstatement}{Definition}{12}\label{def:3.12}
The cross ratio\index{Cross ratio|(}\index{Cross ratio} $(z_{1}, z_{2}, z_{3}, z_{4})$ is the image of $z_{1}$ under the linear transformation which carries $z_{2}, z_{3}, z_{4}$ into 1, 0, $\infty$.
\end{originalstatement}
The definition is meaningful only if $z_{2}, z_{3}, z_{4}$ are distinct. A conventional value can be introduced as soon as any three of the points are distinct, but this is unimportant.


% Original PDF physical page 95
\index{Linear group|)}\sourcepage{95}
The cross ratio\index{Cross ratio} is invariant under linear transformations\index{Linear transformation}. In more precise formulation:

\begin{originalstatement}{Theorem}{12}\label{thm:3.12}
If $z_1, z_2, z_3, z_4$ are distinct points in the extended plane and $T$ any linear transformation, then $(Tz_1, Tz_2, Tz_3, Tz_4) = (z_1, z_2, z_3, z_4)$.
\end{originalstatement}
The proof is immediate, for if $Sz = (z, z_2, z_3, z_4)$, then $ST^{-1}$ carries $Tz_2$, $Tz_3$, $Tz_4$ into $1, 0, \infty$. By definition we have hence
\begin{equation*}
(T z _ {1}, T z _ {2}, T z _ {3}, T z _ {4}) = S T ^ {- 1} (T z _ {1}) = S z _ {1} = (z _ {1}, z _ {2}, z _ {3}, z _ {4}).
\end{equation*}
With the help of this property we can immediately write down the linear transformation which carries three given points $z_{1}$, $z_{2}$, $z_{3}$ to prescribed positions $w_{1}$, $w_{2}$, $w_{3}$. The correspondence must indeed be given by
\begin{equation*}
(w, w _ {1}, w _ {2}, w _ {3}) = (z, z _ {1}, z _ {2}, z _ {3}).
\end{equation*}
In general it is of course necessary to solve this equation with respect to w.

\begin{originalstatement}{Theorem}{13}\label{thm:3.13}
The cross ratio $(z_{1}, z_{2}, z_{3}, z_{4})$ is real if and only if the four points lie on a circle or on a straight line.
\end{originalstatement}
This is evident by elementary geometry, for we obtain
\begin{equation*}
\arg \left(z _ {1}, z _ {2}, z _ {3}, z _ {4}\right) = \arg \frac {z _ {1} - z _ {3}}{z _ {1} - z _ {4}} - \arg \frac {z _ {2} - z _ {3}}{z _ {2} - z _ {4}},
\end{equation*}
and if the points lie on a circle this difference of angles is either 0 or $\pm\pi$, depending on the relative location.

For an analytic proof we need only show that the image of the real axis under any linear transformation is either a circle or a straight line. Indeed, $Tz = (z, z_{2}, z_{3}, z_{4})$ is real on the image of the real axis under the transformation $T^{-1}$ and nowhere else.

The values of $w = T^{-1}z$ for real $z$ satisfy the equation $Tw = \overline{Tw}$. Explicitly, this condition is of the form
\begin{equation*}
\frac {a w + b}{c w + d} = \frac {\bar {a} \bar {w} + \bar {b}}{\bar {c} \bar {w} + \bar {d}}.
\end{equation*}
By cross multiplication we obtain
\begin{equation*}
(a \bar {c} - c \bar {a}) | w ^ {2} | + (a \bar {d} - c \bar {b}) w + (b \bar {c} - d \bar {a}) \bar {w} + b \bar {d} - d \bar {b} = 0.
\end{equation*}
If $a\bar{c} - c\bar{a} = 0$ this is the equation of a straight line, for under this condition the coefficient $a\bar{d} - c\bar{b}$ cannot also vanish. If $a\bar{c} - c\bar{a} \neq 0$ we can
% Original PDF physical page 96
\sourcepage{96}
divide by this coefficient and complete the square. After a simple computation we obtain
\begin{equation*}
\left| w + \frac {\bar {a} d - \bar {c} b}{\bar {a} c - \bar {c} a} \right| = \left| \frac {a d - b c}{\bar {a} c - \bar {c} a} \right|
\end{equation*}
which is the equation of a circle.

The last result makes it clear that we should not, in the theory of linear transformations\index{Linear transformation}, distinguish between circles and straight lines. A further justification was found in the fact that both correspond to circles on the Riemann sphere. Accordingly we shall agree to use the word circle in this wider sense.\hyperlink{fn:96:dagger}{\textsuperscript{\ensuremath{\dagger}}}

The following is an immediate corollary of Theorems~\hyperref[thm:3.12]{12} and \hyperref[thm:3.13]{13}:

\begin{originalstatement}{Theorem}{14}\label{thm:3.14}
A linear transformation carries circles into circles.
\end{originalstatement}
\subsection*{EXERCISES}
\originalnumber{1}{exercise:3.3.2.7.1}
1. Find the linear transformation which carries 0, i, -i into 1, -1, 0.

\originalnumber{2}{exercise:3.3.2.7.2}
2. Express the cross ratios corresponding to the 24 permutations of four points in terms of $\lambda = (z_{1}, z_{2}, z_{3}, z_{4})$.

\originalnumber{3}{exercise:3.3.2.7.3}
3. If the consecutive vertices $z_{1}, z_{2}, z_{3}, z_{4}$ of a quadrilateral lie on a circle, prove that
\begin{equation*}
\left| z _ {1} - z _ {3} \right| \cdot \left| z _ {2} - z _ {4} \right| = \left| z _ {1} - z _ {2} \right| \cdot \left| z _ {3} - z _ {4} \right| + \left| z _ {2} - z _ {3} \right| \cdot \left| z _ {1} - z _ {4} \right|
\end{equation*}
and interpret the result geometrically.

\originalnumber{4}{exercise:3.3.2.7.4}
4. Show that any four distinct points can be carried by a linear transformation to positions 1, -1, k, -k, where the value of k depends on the points. How many solutions are there, and how are they related?

\setcounter{section}{3}\setcounter{subsection}{2}\subsection{Symmetry}\label{sec:3.3.3}
The points z and $\bar{z}$ are symmetric with respect to the real axis. A linear transformation with real coefficients carries the real axis into itself and z, $\bar{z}$ into points which are again symmetric. More generally, if a linear transformation T carries the real axis into a circle C, we shall say that the points w = Tz and $w^{*} = T\bar{z}$ are symmetric with respect to C. This is a relation between w, $w^{*}$ and C which does not depend on T. For if S is another transformation which carries the real axis into C, then $S^{-1}T$ is a real transformation, and hence $S^{-1}w = S^{-1}Tz$ and $S^{-1}w^{*} = S^{-1}T\bar{z}$ are also conjugate. Symmetry\index{Symmetry|(}\index{Symmetry} can thus be defined in the following terms:

\begingroup\renewcommand{\thefootnote}{}\footnotetext{\hypertarget{fn:96:dagger}{}\textsuperscript{\ensuremath{\dagger}} This agreement will be in force only when dealing with linear transformations.}\endgroup

% Original PDF physical page 97
\index{Cross ratio|)}\sourcepage{97}
\begin{originalstatement}{Definition}{13}\index{Distance|(}\label{def:3.13}
The points z and $z^{*}$ are said to be symmetric with respect to the circle C through $z_{1}, z_{2}, z_{3}$ if and only if $(z^{*}, z_{1}, z_{2}, z_{3}) = \overline{(z, z_{1}, z_{2}, z_{3})}$.
\end{originalstatement}
The points on C, and only those, are symmetric to themselves. The mapping which carries z into $z^{*}$ is a one-to-one correspondence and is called reflection with respect to C. Two reflections will evidently result in a linear transformation\index{Linear transformation}.

We wish to investigate the geometric significance of symmetry\index{Symmetry}. Suppose first that C is a straight line. Then we can choose $z_{3} = \infty$ and the condition for symmetry becomes
\begin{equation*}
\frac {z ^ {*} - z _ {2}}{z _ {1} - z _ {2}} = \frac {\bar {z} - \bar {z} _ {2}}{\bar {z} _ {1} - \bar {z} _ {2}}\tag{10}\label{eq:3.10}
\end{equation*}
Taking absolute values we obtain $|z^{*}-z_{2}|=|z-z_{2}|$. Here $z_{2}$ can be any finite point on C, and we conclude that z and $z^{*}$ are equidistant from all points on C. By \hyperref[eq:3.10]{(10)} we have further
\begin{equation*}
\operatorname{Im} \frac {z ^ {*} - z _ {2}}{z _ {1} - z _ {2}} = - \operatorname{Im} \frac {z - z _ {2}}{z _ {1} - z _ {2}},
\end{equation*}
and hence $z$ and $z^*$ are in different half planes determined by $C.$\hyperlink{fn:97:dagger}{\textsuperscript{\ensuremath{\dagger}}}. We leave to the reader to prove that $C$ is the bisecting normal of the segment between $z$ and $z^*$.

Consider now the case of a finite circle C of center a and radius R. Systematic use of the invariance of the cross ratio allows us to conclude as follows:
\begin{equation*}
\begin{aligned}
\overline{(z,z_1,z_2,z_3)}
&=\overline{(z-a,z_1-a,z_2-a,z_3-a)}\\
&=\left(\bar z-\bar a,\frac{R^2}{z_1-a},\frac{R^2}{z_2-a},\frac{R^2}{z_3-a}\right)\\
&=\left(\frac{R^2}{\bar z-\bar a},z_1-a,z_2-a,z_3-a\right)\\
&=\left(\frac{R^2}{\bar z-\bar a}+a,z_1,z_2,z_3\right).
\end{aligned}
\end{equation*}
This equation shows that the symmetric point of $z$ is $z^{*} = R^{2} / (\bar{z} - \bar{a}) + a$ or that $z$ and $z^{*}$ satisfy the relation
\begin{equation*}
(z ^ {*} - a) (z - \bar {a}) = R ^ {2}.\tag{11}\label{eq:3.11}
\end{equation*}
The product $|z^{*}-a|\cdot|z-a|$ of the distances to the center is hence $R^{2}$. Further, the ratio $(z^{*}-a)/(z-a)$ is positive, which means that z and $z^{*}$ are situated on the same half line from a. There is a simple geometric construction for the symmetric point of z (Fig.~\hyperref[fig:3-2]{3-2}). We note that the symmetric point of a is $\infty$.

\begingroup\renewcommand{\thefootnote}{}\footnotetext{\hypertarget{fn:97:dagger}{}\textsuperscript{\ensuremath{\dagger}} Unless they coincide and lie on $C$.}\endgroup

% Original PDF physical page 98
\index{Distance|)}\sourcepage{98}
\bookfigure{98}{1}{FIG. 3-2. Reflection in a circle.}
\begin{originalstatement}{Theorem}{15}\label{thm:3.15}
(The symmetry\index{Symmetry} principle\index{Symmetry principle}.) If a linear transformation\index{Linear transformation} carries a circle $C_{1}$ into a circle $C_{2}$, then it transforms any pair of symmetric points with respect to $C_{1}$ into a pair of symmetric points with respect to $C_{2}$.
\end{originalstatement}
Briefly, linear transformations preserve symmetry. If $C_{1}$ or $C_{2}$ is the real axis, the principle follows from the definition of symmetry. In the general case the assertion follows by use of an intermediate transformation which carries $C_{1}$ into the real axis.

There are two ways in which the principle of symmetry can be used. If the images of z and C under a certain linear transformation are known, then the principle allows us to find the image of $z^{*}$. On the other hand, if the images of z and $z^{*}$ are known, we conclude that the image of C must be a line of symmetry of these images. While this is not enough to determine the image of C, the information we gain is nevertheless valuable.

The principle of symmetry is put to practical use in the problem of finding the linear transformations which carry a circle C into a circle $C'$. We can always determine the transformation by requiring that three points $z_{1}, z_{2}, z_{3}$ on C go over into three points $w_{1}, w_{2}, w_{3}$ on $C'$; the transformation is then $(w, w_{1}, w_{2}, w_{3}) = (z, z_{1}, z_{2}, z_{3})$. But the transformation is also determined if we prescribe that a point $z_{1}$ on C shall correspond to a point $w_{1}$ on $C'$ and that a point $z_{2}$ not on C shall be carried into a point $w_{2}$ not on $C'$. We know then that $z_{2}^{*}$ (the symmetric point of $z_{2}$ with respect to C) must correspond to $w_{2}^{*}$ (the symmetric point of $w_{2}$ with respect to $C'$ ). Hence the transformation will be obtained from the relation $(w, w_{1}, w_{2}, w_{2}^{*}) = (z, z_{1}, z_{2}, z_{2}^{*})$.

\subsection*{EXERCISES}
\originalnumber{1}{exercise:3.3.3.8.1}
1. Prove that every reflection carries circles into circles.


% Original PDF physical page 99
\sourcepage{99}
\originalnumber{2}{exercise:3.3.3.8.2}
2. Reflect the imaginary axis, the line $x = y$, and the circle $|z| = 1$ in the circle $|z - 2| = 1$.

\originalnumber{3}{exercise:3.3.3.8.3}
3. Carry out the reflections in the preceding exercise by geometric construction.

\originalnumber{4}{exercise:3.3.3.8.4}
4. Find the linear transformation\index{Linear transformation} which carries the circle $|z| = 2$ into $|z + 1| = 1$, the point -2 into the origin, and the origin into i.

\originalnumber{5}{exercise:3.3.3.8.5}
5. Find the most general linear transformation of the circle $|z| = R$ into itself.

\originalnumber{6}{exercise:3.3.3.8.6}
6. Suppose that a linear transformation carries one pair of concentric circles into another pair of concentric circles. Prove that the ratios of the radii must be the same.

\originalnumber{7}{exercise:3.3.3.8.7}
7. Find a linear transformation which carries $|z| = 1$ and $|z - \frac{1}{4}| = \frac{1}{4}$ into concentric circles. What is the ratio of the radii?

\originalnumber{8}{exercise:3.3.3.8.8}
8. Same problem for $|z| = 1$ and $x = 2$.

\setcounter{section}{3}\setcounter{subsection}{3}\subsection{Oriented Circles}\label{sec:3.3.4}
Because $S(z)$ is analytic and
\begin{equation*}
S ^ {\prime} (z) = \frac {a d - b c}{(c z + d) ^ {2}} \neq 0
\end{equation*}
the mapping $w = S(z)$ is conformal for $z \neq -d/c$ and $\infty$. It follows that a pair of intersecting circles are mapped on circles that include the same angle. In addition, the sense of an angle is preserved. From an intuitive point of view this means that right and left are preserved, but a more precise formulation is desirable.

An orientation\index{Orientation} of a circle C is determined by an ordered triple of points $z_{1}, z_{2}, z_{3}$ on C. With respect to this orientation a point z not on C is said to lie to the right of C if $\operatorname{Im}(z, z_{1}, z_{2}, z_{3}) > 0$ and to the left of C is $\operatorname{Im}(z, z_{1}, z_{2}, z_{3}) < 0$ (this checks with everyday use because $(i, 1, 0, \infty) = i$ ). It is essential to show that there are only two different orientations. By this we mean that the distinction between left and right is the same for all triples, while the meaning may be reversed. Since the cross ratio is invariant, it is sufficient to consider the case where C is the real axis. Then
\begin{equation*}
(z, z _ {1}, z _ {2}, z _ {3}) = \frac {a z + b}{c z + d}
\end{equation*}
can be written with real coefficients, and a simple calculation gives
\begin{equation*}
\operatorname{Im} \left(z, z _ {1}, z _ {2}, z _ {3}\right) = \frac {a d - b c}{| c z + d | ^ {2}} \operatorname{Im} z.
\end{equation*}
We recognize that the distinction between right and left is the same as the distinction between the upper and lower half plane. Which is which depends on the sign of the determinant ad - bc.


% Original PDF physical page 100
\index{Symmetry|)}\sourcepage{100}
A linear transformation\index{Linear transformation} S carries the oriented circle C into a circle which we orient through the triple $Sz_{1}$, $Sz_{2}$, $Sz_{3}$. From the invariance of the cross ratio it follows that the left and right of C will be mapped on the left and right of the image circle.

If two circles are tangent to each other, their orientations can be compared. Indeed, we can use a linear transformation which throws their common point to $\infty$. The circles become parallel straight lines, and we know how to compare the directions of parallel lines.

In the geometric representation the orientation $z_{1}, z_{2}, z_{3}$ can be indicated by an arrow which points from $z_{1}$ over $z_{2}$ to $z_{3}$. With the usual choice of the coordinate system left and right will have their customary meaning with respect to this arrow.

When the finite plane is considered as part of the extended plane, the point at infinity is distinguished. We can therefore define an absolute positive orientation of all finite circles by the requirement that $\infty$ should lie to the right of the oriented circles. The points to the left are said to form the inside of the circle and the points to the right form its outside.

\subsection*{EXERCISES}
\originalnumber{1}{exercise:3.3.4.9.1}
1. If $z_{1}, z_{2}, z_{3}, z_{4}$ are points on a circle, show that $z_{1}, z_{3}, z_{4}$ and $z_{2}, z_{3}, z_{4}$ determine the same orientation if and only if $(z_{1}, z_{2}, z_{3}, z_{4}) > 0$.

\originalnumber{2}{exercise:3.3.4.9.2}
2. Prove that a tangent to a circle is perpendicular to the radius through the point of contact (in this connection a tangent should be defined as a straight line with only one point in common with the circle).

\originalnumber{3}{exercise:3.3.4.9.3}
3. Verify that the inside of the circle $|z - a| = R$ is formed by all points z with $|z - a| < R$.

\originalnumber{4}{exercise:3.3.4.9.4}
4. The angle\index{Angle} between two oriented circles at a point of intersection is defined as the angle between the tangents at that point, equipped with the same orientation. Prove by analytic reasoning, rather than geometric inspection, that the angles at the two points of intersection are opposite to each other.

\setcounter{section}{3}\setcounter{subsection}{4}\subsection{Families of Circles}\label{sec:3.3.5}
A great deal can be done toward the visualization of linear transformations by the introduction of certain families of circles which may be thought of as coordinate lines in a circular coordinate system.

Consider a linear transformation of the form
\begin{equation*}
w = k \cdot \frac {z - a}{z - b}.
\end{equation*}
Here z = a corresponds to w = 0 and z = b to $w = \infty$. It follows that the straight lines through the origin of the w-plane are images of the
% Original PDF physical page 101
\sourcepage{101}
circles through $a$ and $b$. On the other hand, the concentric circles about the origin, $|w| = \rho$, correspond to circles with the equation
\begin{equation*}
\left| \frac {z - a}{z - b} \right| = \rho / | k |.
\end{equation*}
These are the circles of Apollonius\index{Apollonius} with limit points a and b. By their equation they are the loci of points whose distances from a and b have a constant ratio.

Denote by $C_{1}$ the circles through a, b and by $C_{2}$ the circles of Apollonius with these limit points. The configuration (Fig.~\hyperref[fig:3-3]{3-3}) formed by all the circles $C_{1}$ and $C_{2}$ will be referred to as the circular net or the Steiner\index{Steiner J@Steiner, J.} circles determined by a and b. It has many interesting properties of which we shall list a few:

1. There is exactly one $C_{1}$ and one $C_{2}$ through each point in the plane with the exception of the limit points.

2. Every $C_1$ meets every $C_2$ under right angles.

3. Reflection in a $C_{1}$ transforms every $C_{2}$ into itself and every $C_{1}$ into another $C_{1}$. Reflection in a $C_{2}$ transforms every $C_{1}$ into itself and every $C_{2}$ into another $C_{2}$.

4. The limit points are symmetric with respect to each $C_{2}$, but not with respect to any other circle.

\bookfigure{101}{1}{FIG. 3-3. Steiner circles.}

% Original PDF physical page 102
\sourcepage{102}
These\index{Parabolic transformation|(} properties are all trivial when the limit points are 0 and $\infty$, i.e., when the $C_{1}$ are lines through the origin and the $C_{2}$ concentric circles. Since the properties are invariant under linear transformations\index{Linear transformation}, they must continue to hold in the general case.

If a transformation $w = Tz$ carries $a, b$ into $a'$, $b'$ it can be written in the form
\begin{equation*}
\frac {w - a ^ {\prime}}{w - b ^ {\prime}} = k \cdot \frac {z - a}{z - b}.\tag{12}\label{eq:3.12}
\end{equation*}
It is clear that $T$ transforms the circles $C_1$ and $C_2$ into circles $C_1'$ and $C_2'$ with the limit points $a', b'$.

The situation is particularly simple if $a' = a$, $b' = b$. Then a, b are said to be fixed points\index{Fixed point}\index{Point!fixed} of T, and it is convenient to represent z and Tz in the same plane. Under these circumstances the whole circular net will be mapped upon itself. The value of k serves to identify the image circles $C_{1}'$ and $C_{2}'$. Indeed, with appropriate orientations $C_{1}$ forms the angle arg k with its image $C_{1}'$, and the quotient of the constant ratios $|z - a|/|z - b|$ on $C_{2}'$ and $C_{2}$ is $|k|$.

The special cases in which all $C_{1}$ or all $C_{2}$ are mapped upon themselves are particularly important. We have $C_{1}^{\prime}=C_{1}$ for all $C_{1}$ if k>0 (if k<0 the circles are still the same, but the orientation is reversed). The transformation is then said to be hyperbolic. When k increases the points $Tz, z\neq a, b$, will flow along the circles $C_{1}$ toward b. The consideration of this flow provides a very clear picture of a hyperbolic transformation\index{Hyperbolic transformation}.

The case $C_{2}^{\prime}=C_{2}$ occurs when $|k|=1$. Transformations with this property are called elliptic. When arg k varies, the points Tz move along the circles $C_{2}$. The corresponding flow circulates about a and b in different directions.

The general linear transformation with two fixed points is the product of a hyperbolic and an elliptic transformation\index{Elliptic transformation} with the same fixed points.

The fixed points of a linear transformation are found by solving the equation
\begin{equation*}
z = \frac {\alpha z + \beta}{\gamma z + \delta}.\tag{13}\label{eq:3.13}
\end{equation*}
In general this is a quadratic equation with two roots; if $\gamma = 0$ one of the fixed points is $\infty$. It may happen, however, that the roots coincide. A linear transformation with coinciding fixed points is said to be parabolic. The condition for this is $(\alpha - \delta)^{2} = 4\beta\gamma$.

If the equation \hyperref[eq:3.13]{(13)} is found to have two distinct roots $a$ and $b$, the transformation can be written in the form


% Original PDF physical page 103
\index{Parabolic transformation|)}\sourcepage{103}
\begin{equation*}
{\frac {w - a}{w - b}} = k {\frac {z - a}{z - b}}.
\end{equation*}
We can then use the Steiner circles determined by a, b to discuss the nature of the transformation. It is important to note, however, that the method is by no means restricted to this case. We can write any linear transformation\index{Linear transformation} in the form \hyperref[eq:3.12]{(12)} with arbitrary a, b and use the two circular nets to great advantage.

For the discussion of parabolic transformations it is desirable to introduce still another type of circular net. Consider the transformation
\begin{equation*}
w = \frac {\omega}{z - a} + c.
\end{equation*}
It is evident that straight lines in the w-plane correspond to circles through a; moreover, parallel lines correspond to mutually tangent circles. In particular, if $w = u + iv$ the lines u = constant and v = constant correspond to two families of mutually tangent circles which intersect at right angles (Fig.~\hyperref[fig:3-4]{3-4}). This configuration can be considered as a degenerate set of Steiner circles. It is determined by the point a and the tangent to one of the families of circles. We shall denote the images of the lines v = constant by $C_{1}$, the circles of the other family by $C_{2}$. Clearly, the line v = Im c corresponds to the tangent of the circles $C_{1}$; its direction is given by arg $\omega$.

\bookfigure{103}{1}{FIG. 3-4. Degenerate Steiner circles.}

% Original PDF physical page 104
\sourcepage{104}
Any\index{Loxodromic transformation|(} transformation which carries a into $a'$ can be written in the form
\begin{equation*}
\frac {\omega^ {\prime}}{w - a ^ {\prime}} = \frac {\omega}{z - a} + c.
\end{equation*}
It is clear that the circles $C_{1}$ and $C_{2}$ are carried into the circles $C_{1}^{\prime}$ and $C_{2}^{\prime}$ determined by $a^{\prime}$ and $\omega^{\prime}$. We suppose now that $a = a^{\prime}$ is the only fixed point. Then $\omega = \omega^{\prime}$ and we can write
\begin{equation*}
\frac {\omega}{w - a} = \frac {\omega}{z - a} + c.\tag{14}\label{eq:3.14}
\end{equation*}
By this transformation the configuration consisting of the circles $C_{1}$ and $C_{2}$ is mapped upon itself. In \hyperref[eq:3.14]{(14)} a multiplicative factor is arbitrary, and we can hence suppose that c is real. Then every $C_{1}$ is mapped upon itself and the parabolic transformation can be considered as a flow along the circles $C_{2}$.

A linear transformation\index{Linear transformation} that is neither hyperbolic, elliptic, nor parabolic is said to be loxodromic.

\subsection*{EXERCISES}
\originalnumber{1}{exercise:3.3.5.10.1}
1. Find the fixed points of the linear transformations
\begin{equation*}
w = \frac {z}{2 z - 1}, \quad w = \frac {2 z}{3 z - 1}, \quad w = \frac {3 z - 4}{z - 1}, \quad w = \frac {z}{2 - z}.
\end{equation*}
Is any of these transformations elliptic, hyperbolic, or parabolic?

\originalnumber{2}{exercise:3.3.5.10.2}
2. Suppose that the coefficients of the transformation
\begin{equation*}
S z = \frac {a z + b}{c z + d}
\end{equation*}
are normalized by $ad - bc = 1$. Show that $S$ is elliptic if and only if $-2 < a + d < 2$, parabolic if $a + d = \pm 2$, hyperbolic if $a + d < -2$ or $>2$.

\originalnumber{3}{exercise:3.3.5.10.3}
3. Show that a linear transformation which satisfies $S^{n}z = z$ for some integer n is necessarily elliptic.

\originalnumber{4}{exercise:3.3.5.10.4}
4. If S is hyperbolic or loxodromic, show that $S^{n}z$ converges to a fixed point as $n \to \infty$, the same for all z, except when z coincides with the other fixed point. (The limit is the attractive, the other the repellent fixed point. What happens when $n \to -\infty$? What happens in the parabolic case?)

\originalnumber{5}{exercise:3.3.5.10.5}
5. Find all linear transformations which represent rotations of the Riemann sphere.

\originalnumber{6}{exercise:3.3.5.10.6}
6. Find all circles which are orthogonal to $|z| = 1$ and $|z - 1| = 4$.


% Original PDF physical page 105
\index{Loxodromic transformation|)}\sourcepage{105}
\originalnumber{7}{exercise:3.3.5.10.7}
7. In an obvious way, which we shall not try to make precise, a family of transformations depends on a certain number of real parameters. How many real parameters are there in the family of all linear transformations\index{Linear transformation}? How many in the families of hyperbolic, elliptic, parabolic transformations? How many linear transformations leave a given circle C invariant?

\setcounter{section}{3}\section{Elementary Conformal Mappings}\label{sec:3.4}
The conformal mapping\index{Conformal mapping} associated with an analytic function affords an excellent visualization of the properties of the latter; it can well be compared with the visualization of a real function by its graph. It is therefore natural that all questions connected with conformal mapping have received a great deal of attention; progress in this direction has increased our knowledge of analytic functions considerably. In addition, conformal mapping enters naturally in many branches of mathematical physics and in this way accounts for the immediate usefulness of complex-function theory.

One of the most important problems is to determine the conformal mappings of one region onto another. In this section we shall consider those mappings which can be defined by elementary functions.

\setcounter{section}{4}\setcounter{subsection}{0}\subsection{The Use of Level Curves}\label{sec:3.4.1}
When a conformal mapping is defined by an explicit analytic function $w = f(z)$, we naturally wish to gain information about the specific geometric properties of the mapping. One of the most fruitful ways is to study the correspondence of curves induced by the point transformation. The special properties of the function $f(z)$ may express themselves in the fact that certain simple curves are transformed into curves of a family of well-known character. Any such information will strengthen our visual conception of the mapping.

Such was the case for mappings by linear transformations. We proved in Sec. 3 that a linear transformation carries circles into circles, provided that straight lines are included as a special case. By consideration of the Steiner circles it was possible to obtain a complete picture of the correspondence.

In more general cases it is advisable to begin with a study of the image curves of the lines $x = x_{0}$ and $y = y_{0}$. If we write $f(z) = u(x,y) + iv(x,y)$, the image of $x = x_{0}$ is given by the parametric equations $u = u(x_{0},y)$, $v = v(x_{0},y)$; y acts as a parameter and can be eliminated or retained according to convenience. The image of $y = y_{0}$ is determined in the same way. Together, the curves form an orthogonal net in the w-plane. Similarly, we may consider the curves $u(x,y) = u_{0}$ and $v(x,y) = v_{0}$ in the z-plane. They are also orthogonal and are called the level curves\index{Curve!level}\index{Level curve} of u and v.


% Original PDF physical page 106
\index{Linear transformation|)}\sourcepage{106}
In\index{Parabola|(} other cases it may be more convenient to use polar coordinates and study the images of concentric circles and straight lines through the origin.

Among the simplest mappings are those by a power $w = z^{\alpha}$. We consider only the case of real $\alpha$, and then we may as well suppose that $\alpha$ is positive. Since
\begin{equation*}
\begin{array}{r c l} | w | & = & | z | ^ {\alpha} \\ \arg w & = & \alpha \arg z \end{array}
\end{equation*}
concentric circles about the origin are transformed into circles of the same family, and half lines from the origin correspond to other half lines. The mapping is conformal at all points $z \neq 0$, but an angle $\theta$ at the origin is transformed into an angle $\alpha\theta$. For $\alpha \neq 1$ the transformation of the whole plane is not one to one, and if $\alpha$ is fractional $z^{\alpha}$ is not even single-valued. In general we can therefore only consider the mapping of an angular sector onto another.

The sector $S(\varphi_1, \varphi_2)$, where $0 < \varphi_2 - \varphi_1 \leq 2\pi$, shall be formed by all points $z \neq 0$ such that one value of $\arg z$ satisfies the inequality
\begin{equation*}
\varphi_ {1} <   \arg z <   \varphi_ {2}.\tag{15}\label{eq:3.15}
\end{equation*}
It is easy to show that $S(\varphi_1, \varphi_2)$ is a region. In this region a unique value of $w = z^\alpha$ is defined by the condition
\begin{equation*}
\arg w = \alpha \arg z
\end{equation*}
where $\arg z$ stands for the value of the argument singled out by the condition \hyperref[eq:3.15]{(15)}. This function is analytic with the nonvanishing derivative
\begin{equation*}
D e ^ {\alpha \log z} = \alpha \frac {w}{z}.
\end{equation*}
The mapping is one to one only if $\alpha(\varphi_{2}-\varphi_{1})\leq2\pi$, and in this case $S(\varphi_{1},\varphi_{2})$ is mapped onto the sector $S(\alpha\varphi_{1},\alpha\varphi_{2})$ in the w-plane. It should be observed that $S(\varphi_{1}+n\cdot2\pi,\varphi_{2}+n\cdot2\pi)$ is geometrically identical with $S(\varphi_{1},\varphi_{2})$ but may determine a different branch of $z^{\alpha}$.

Let us consider the mapping $w = z^{2}$ in detail. Since $u = x^{2} - y^{2}$ and v = 2xy, we recognize that the level curves $u = u_{0}$ and $v = v_{0}$ are equilateral hyperbolas with the diagonals and the coordinate axes for asymptotes. They are of course orthogonal to each other. On the other hand, the image of $x = x_{0}$ is $v^{2} = 4x_{0}^{2}(x_{0}^{2} - u)$ and the image of $y = y_{0}$ is $v^{2} = 4y_{0}^{2}(y_{0}^{2} + u)$. Both families represent parabolas with the focus at the origin whose axes are pointed in the negative and positive direction of the u-axis. Their orthogonality is well-known from analytic geometry. The families of level curves are shown in Figs. 3-5 and 3-6.


% Original PDF physical page 107
\index{Parabola|)}\sourcepage{107}
\bookfigure{107}{1}{FIG. 3-5. z-plane.}
For a different family of image curves consider the circles $|w - 1| = k$ in the w-plane. The equation of the inverse image can be written in the form
\begin{equation*}
(x ^ {2} + y ^ {2}) ^ {2} = 2 (x ^ {2} - y ^ {2}) + k ^ {2} - 1
\end{equation*}
and represents a family of lemniscates with the focal points $\pm1$. The orthogonal family is represented by
\begin{equation*}
x ^ {2} - y ^ {2} = 2 h x y + 1
\end{equation*}
and consists of all equilateral hyperbolas with center at the origin which pass through the points $\pm1$.

In the case of the third power $w = z^{3}$ the level curves in both planes are cubic curves. There is no point in deriving their equations, for their general shape is clear without calculation. For instance, the curves $u = u_{0} > 0$ must have the form indicated in Fig.~\hyperref[fig:3-7]{3-7}. Similarly, if we follow the change of arg w when z traces the line $x = x_{0} > 0$, we find that the image curve must have a loop (Fig.~\hyperref[fig:3-8]{3-8}). It is a folium of Descartes.


% Original PDF physical page 108
\sourcepage{108}
\bookfigure{108}{1}{FIG. 3-6. w-plane.}
The mapping by $w = e^{z}$ is very simple. The lines $x = x_{0}$ and $y = y_{0}$ are mapped onto circles about the origin and rays of constant argument. Any other straight line in the z-plane is mapped on a logarithmic spiral. The mapping is one to one in any region which does not contain two points whose difference is a multiple of $2\pi i$. In particular, a horizontal strip $y_{1} < y < y_{2}$, $y_{2} - y_{1} \leq 2\pi$ is mapped onto an angular sector, and if $y_{2} - y_{1} = \pi$ the image is a half plane. We are thus able to map a parallel strip onto a half plane, and hence onto any circular region. The left half of the strip, cut off by the imaginary axis, corresponds to a half circle.

It is useful to write down some explicit formulas for the mapping. The function $\zeta = \xi + i\eta = e^{z}$ maps the strip $-\pi/2 < y < \pi/2$ onto the half plane $\xi > 0$. On the other hand,
\begin{equation*}
w = \frac {\zeta - 1}{\zeta + 1}
\end{equation*}
% Original PDF physical page 109
\sourcepage{109}
maps $\xi > 0$ onto $|w| < 1$. Hence
\begin{equation*}
w = \frac {e ^ {z} - 1}{e ^ {z} + 1} = \tanh \frac {z}{2}
\end{equation*}
maps the strip $|\operatorname{Im} z| < \pi/2$ on the unit disk $|w| < 1$.

\setcounter{section}{4}\setcounter{subsection}{1}\subsection{A Survey of Elementary Mappings}\label{sec:3.4.2}
When faced with the problem of mapping a region $\Omega_{1}$ conformally onto another region $\Omega_{2}$, it is usually advisable to proceed in two steps. First, we map $\Omega_{1}$ onto a circular region, and then we map the circular region onto $\Omega_{2}$. In other words, the general problem of conformal mapping\index{Conformal mapping} can be reduced to the problem of mapping a region onto a disk or a half plane. We shall prove, in Chap.~\hyperref[chap:6]{6}, that this mapping problem has a solution for every region whose boundary consists of a simple closed curve.

The main tools at our disposal are linear transformations and transformations by a power, by the exponential function, and by the logarithm. All these transformations have the characteristic property that they map a family of straight lines or circles on a similar family. For this reason, their use is essentially limited to regions whose boundary is made up of circular arcs and line segments. The power serves the particular purpose of straightening angles, and with the aid of the exponential function we can even transform zero angles into straight angles.

By these means we can for instance find a standard mapping of any
\bookfigure{109}{1}{FIG. 3-7}
\bookfigure{109}{2}{FIG. 3-8}
% Original PDF physical page 110
\sourcepage{110}
region whose boundary consists of two circular arcs with common end points. Such a region is either a circular wedge, whose angle may be greater than $\pi$, or its complement. If the end points of the arcs are a and b, we begin with the preliminary mapping $z_{1} = (z - a)/(z - b)$ which transforms the given region into an angular sector. By an appropriate power $w = z_{1}^{\alpha}$ this sector can be mapped onto a half plane.

If the circles are tangent to each other at the point a, the transformation $z_{1}=1/(z-a)$ will map the region between them onto a parallel strip, and a suitable exponential transformation maps the strip onto a half plane.

A little more generally, the same method applies to a circular triangle with two right angles. In fact, if the third angle has the vertex $a$, and if the sides from $a$ meet again at $b$, the linear transformation $z_{1} = (z - a)/(z - b)$ maps the triangle onto a circular sector. By means of a power this sector can be transformed into a half circle; the half circle is a wedge-shaped region which in turn can be mapped onto a half plane.

In this connection we shall treat explicitly a special case which occurs frequently. Let it be required to map the complement of a line segment onto the inside or outside of a circle. The region is a wedge with the angle $2\pi$; without loss of generality we may assume that the end points of the segment are $\pm1$. The preliminary transformation
\begin{equation*}
z _ {1} = \frac {z + 1}{z - 1}
\end{equation*}
maps the wedge on the full angle obtained by exclusion of the negative real axis. Next we define
\begin{equation*}
z _ {2} = \sqrt {z _ {1}}
\end{equation*}
as the square root whose real part is positive and obtain a map onto the right half plane. The final transformation
\begin{equation*}
w = \frac {z _ {2} - 1}{z _ {2} + 1}
\end{equation*}
maps the half plane onto $|w| < 1$.

Elimination of the intermediate variables leads to the correspondence
\begin{equation*}
\begin{array}{c} z = \frac {1}{2} \left(w + \frac {1}{w}\right) \\ w = z - \sqrt {z ^ {2} - 1}. \end{array}\tag{16}\label{eq:3.16}
\end{equation*}
The sign of the square root is uniquely determined by the condition $|w| < 1$, for $(z - \sqrt{z^{2} - 1})(z + \sqrt{z^{2} - 1}) = 1$. If the sign is changed, we obtain a mapping onto $|w| > 1$.


% Original PDF physical page 111
\sourcepage{111}
For a more detailed study of the mapping \hyperref[eq:3.16]{(16)} we set $w = \rho e^{i\theta}$ and obtain
\begin{equation*}
x = \frac {1}{2} \left(\rho + \frac {1}{\rho}\right) \cos \theta
\end{equation*}
\begin{equation*}
y = \frac {1}{2} \left(\rho - \frac {1}{\rho}\right) \sin \theta .
\end{equation*}
Elimination of $\theta$ yields
\begin{equation*}
\frac {x ^ {2}}{[ \frac {1}{2} (\rho + \rho^ {- 1}) ] ^ {2}} + \frac {y ^ {2}}{[ \frac {1}{2} (\rho - \rho^ {- 1}) ] ^ {2}} = 1\tag{17}\label{eq:3.17}
\end{equation*}
and elimination of $\rho$
\begin{equation*}
\frac {x ^ {2}}{\cos^ {2} \theta} - \frac {y ^ {2}}{\sin^ {2} \theta} = 1.\tag{18}\label{eq:3.18}
\end{equation*}
Hence the image of a circle $|w| = \rho < 1$ is an ellipse\index{Ellipse} with the major axis $\rho + \rho^{-1}$ and the minor axis $\rho^{-1} - \rho$. The image of a radius is half a branch of a hyperbola\index{Hyperbola|(}\index{Hyperbola}. The ellipses \hyperref[eq:3.17]{(17)} and the hyperbolas \hyperref[eq:3.18]{(18)} are confocal. The correspondence is illustrated in Fig.~\hyperref[fig:3-9]{3-9}.

Clearly, the transformation \hyperref[eq:3.16]{(16)} allows us to include in our list of elementary conformal mappings\index{Conformal mapping} the mapping of the outside of an ellipse or the region between the branches of a hyperbola onto a circular region. It does not, however, allow us to map the inside of an ellipse or the inside of a hyperbolic branch.

As a final and less trivial example we shall study the mapping defined by a cubic polynomial $w = a_{0}z^{3} + a_{1}z^{2} + a_{2}z + a_{3}$. The familiar transformation $z = z_{1} - a_{1}/3a_{0}$ allows us to get rid of the quadratic term,
\bookfigure{111}{1}{FIG. 3-9. Mapping by $z = \frac{1}{2}(w + w^{-1})$.}
% Original PDF physical page 112
\sourcepage{112}
and by obvious normalizations we can reduce the polynomial to the form $w = z^{3} - 3z$. The coefficient for z is chosen so as to make the derivative vanish for $z = \pm 1$.

Making use of the transformation \hyperref[eq:3.16]{(16)} we introduce an auxiliary variable $\zeta$ defined by
\begin{equation*}
z = \zeta + \frac {1}{\zeta}.
\end{equation*}
Our cubic polynomial takes then the simple form
\begin{equation*}
w = \zeta^ {3} + \frac {1}{\zeta^ {3}}.
\end{equation*}
We note that each z determines two values $\zeta$, but they are reciprocal and yield the same value of w. In order to obtain a unique $\zeta$ we may impose the condition $|\zeta| < 1$, but then the segment $(-2,2)$ must be excluded from the z-plane.

It is now easy to visualize the correspondence between the z- and w-planes. To the circle $|\zeta| = \rho < 1$ corresponds an ellipse with the semiaxes $\rho^{-1} \pm \rho$ in the z-plane, and one with the semiaxes $\rho^{-3} \pm \rho^{3}$ in the w-plane. Similarly, a radius $\arg \zeta = \theta$ corresponds to hyperbolic branches in the z- and w-planes; the one in the z-plane has an asymptote which makes the angle $-\theta$ with the positive real axis, and in the w-plane the corresponding angle is $-3\theta$. The whole pattern of confocal ellipses and hyperbolas remains invariant, but when z describes an ellipse w will trace the corresponding larger ellipse three times. The situation is thus very similar to the one in the case of the simpler mapping $w = z^{3}$. For orientation the reader may again lean on Fig.~\hyperref[fig:3-9]{3-9}.

For the region between two hyperbolic branches whose asymptotes make an angle $\leq 2\pi/3$ the mapping is one to one. We note in particular that the six regions into which the hyperbola\index{Hyperbola} $3x^{2} - y^{2} = 3$ and the x-axis divide the z-plane are mapped onto half planes, three of them onto the upper half plane and three onto the lower. The inside of the right-hand branch of the hyperbola corresponds to the whole w-plane with an incision along the negative real axis up to the point -2.

\subsection*{EXERCISES}
All mappings are to be conformal.

\originalnumber{1}{exercise:3.4.2.11.1}
1. Map the common part of the disks $|z| < 1$ and $|z - 1| < 1$ on the inside of the unit circle. Choose the mapping so that the two symmetries are preserved.

\originalnumber{2}{exercise:3.4.2.11.2}
2. Map the region between $|z| = 1$ and $|z - \frac{1}{2}| = \frac{1}{2}$ on a half plane.


% Original PDF physical page 113
\sourcepage{113}
\originalnumber{3}{exercise:3.4.2.11.3}
3.\index{Lindelöf E@Lindelöf, E.|(} Map the complement of the arc $|z| = 1$, $y \geq 0$ on the outside of the unit circle so that the points at $\infty$ correspond to each other.

\originalnumber{4}{exercise:3.4.2.11.4}
4. Map the outside of the parabola $y^2 = 2px$ on the disk $|w| < 1$ so that $z = 0$ and $z = -p/2$ correspond to $w = 1$ and $w = 0$. (Lindelöf.)

\originalnumber{5}{exercise:3.4.2.11.5}
5. Map the inside of the right-hand branch of the hyperbola\index{Hyperbola} $x^{2}-y^{2}=a^{2}$ on the disk $|w|<1$ so that the focus corresponds to w=0 and the vertex to w=-1. (Lindelöf.)

\originalnumber{6}{exercise:3.4.2.11.6}
6. Map the inside of the lemniscate $|z^{2}-a^{2}|=\rho^{2}(\rho>a)$ on the disk $|w|<1$ so that symmetries are preserved. (Lindelöf.)

\originalnumber{7}{exercise:3.4.2.11.7}
7. Map the outside of the ellipse $(x/a)^{2} + (y/b)^{2} = 1$ onto $|w| < 1$ with preservation of symmetries.

\originalnumber{8}{exercise:3.4.2.11.8}
8. Map the part of the z-plane to the left of the right-hand branch of the hyperbola $x^{2} - y^{2} = 1$ on a half plane. (Lindelöf.)

Hint: Consider on one side the mapping of the upper half of the region by $w = z^{2}$, on the other side the mapping of a quadrant by
\begin{equation*}
w = z ^ {3} - 3 z.
\end{equation*}
\setcounter{section}{4}\setcounter{subsection}{2}\subsection{Elementary Riemann Surfaces}\label{sec:3.4.3}
The visualization of a function by means of the corresponding mapping is completely clear only when the mapping is one to one. If this is not the case, we can still give our imagination the necessary support by the introduction of generalized regions in which distinct points may have the same coordinates. In order to do this it is necessary to suppose that points which occupy the same place can be distinguished by other characteristics, for instance a tag or a color. Points with the same tag are considered to lie in the same sheet\index{Sheet} or layer.

This idea leads to the notion of a Riemann surface\index{Riemann surface|(}\index{Riemann surface}. It is not our intention to give, in this connection, a rigorous definition of this notion. For our purposes it is sufficient to introduce Riemann surfaces in a purely descriptive manner. We are free to do so as long as we use them merely for purposes of illustration, and never in logical proofs.

The simplest Riemann surface is connected with the mapping by a power $w = z^{n}$, where n > 1 is an integer. We know that there is a one-to-one correspondence between each angle $(k - 1)(2\pi/n) < \arg z < k(2\pi/n)$, $k = 1, \ldots, n$, and the whole w-plane except for the positive real axis. The image of each angle is thus obtained by performing a “cut” along the positive axis; this cut has an upper and a lower “edge.” Corresponding to the n angles in the z-plane we consider n identical copies of the w-plane with the cut. They will be the “sheets” of the Riemann surface, and they are distinguished by a tag k which serves to identify the corresponding angle. When z moves in its plane, the corresponding
% Original PDF physical page 114
\index{Hyperbola|)}\index{Lindelöf E@Lindelöf, E.|)}\sourcepage{114}
point\index{Order, algebraic!of a branch point|(} w should be free to move on the Riemann surface\index{Riemann surface}. For this reason we must attach the lower edge of the first sheet to the upper edge of the second sheet, the lower edge of the second sheet to the upper edge of the third, and so on. In the last step the lower edge of the nth sheet is attached to the upper edge of the first sheet, completing the cycle. In a physical sense this is not possible without self-intersection, but the idealized model shall be free from this discrepancy. The result of the construction is a Riemann surface whose points are in one-to-one correspondence with the points of the z-plane. What is more, this correspondence is continuous if continuity is defined in the sense suggested by the construction.

The cut along the positive axis could be replaced by a cut along any simple arc from 0 to $\infty$; the Riemann surface obtained in this way should be considered as identical with the one originally constructed. In other words, the cuts are in no way distinguished lines on the surface, but the introduction of specific cuts is necessary for descriptive purposes.

The point w = 0 is in a special position. It connects all the sheets, and a curve must wind n times around the origin before it closes. A point of this kind is called a branch point\index{Branch point}\index{Point!branch}. If our Riemann surface is considered over the extended plane, the point at $\infty$ is also a branch point. In more general cases a branch point need not connect all the sheets; if it connects h sheets, it is said to be of order h - 1.

The Riemann surface corresponding to $w = e^{z}$ is of similar nature. In this case the function maps each parallel strip $(k - 1)2\pi < y < k \cdot 2\pi$ onto a sheet with a cut along the positive axis. The sheets are attached to each other so that they form an endless screw. The origin will not be a point of the Riemann surface, corresponding to the fact that $e^{z}$ is never zero.

The reader will find it easy to construct other Riemann surfaces. We will illustrate the procedure by consideration of the Riemann surface defined by $w = \cos z$. A region which is mapped in a one-to-one manner onto the whole plane, except for one or more cuts, is called a fundamental region\index{Fundamental region|(}\index{Fundamental region}. For fundamental regions of $w = \cos z$ we may choose the strips $(k - 1)\pi < x < k\pi$. Each strip is mapped onto the whole w-plane with cuts along the real axis from $-\infty$ to -1 and from 1 to $\infty$. The line $x = k\pi$ corresponds to both edges of the positive cut if k is even, and

\bookfigure{114}{1}{FIG. 3-10. The Riemann surface of cos z.}

% Original PDF physical page 115
\index{Order, algebraic!of a branch point|)}\sourcepage{115}
\bookfigure{115}{1}{FIG. 3-11. Fundamental regions of cos z.}
to the edges of the negative cut if k is odd. If we consider the two strips which are adjacent along the line $x = k\pi$, we find that the edges of the corresponding cuts must be joined crosswise so as to generate a simple branch point at $w = \pm1$. The resulting surface has infinitely many simple branch points over w = 1 and w = -1 which alternatingly connect the odd and even sheets.

An attempt to illustrate the connection between the sheets is made in Fig.~\hyperref[fig:3-10]{3-10}. It represents a cross section of the surface in the case that the cuts are chosen parallel to each other. The reader should bear in mind that any two points on the same level can be joined by an arc which does not intersect any of the cuts.

Whatever the advantage of such representations may be, the clearest picture of the Riemann surface\index{Riemann surface} is obtained by direct consideration of the fundamental regions\index{Fundamental region} in the z-plane. The interpretation is even simpler if, as in Fig.~\hyperref[fig:3-11]{3-11}, we introduce the subregions which correspond to the upper and lower half plane. The shaded regions are those in which cos z has a positive imaginary part. Each region corresponds to a half plane on which we mark the boundary points 1 and -1. For any two adjacent regions, one white and one shaded, the half planes must be joined across one of the three intervals $(-∞,-1)$, $(-1,1)$, $(1,∞)$. The choice of the correct junction is automatic from a glance at the corresponding situation in the z-plane.

\subsection*{EXERCISES}
\originalnumber{1}{exercise:3.4.3.12.1}
1. Describe the Riemann surface associated with the function
\begin{equation*}
w = \frac {1}{2} \left(z + \frac {1}{z}\right).
\end{equation*}
\originalnumber{2}{exercise:3.4.3.12.2}
2. Same problem for $w = (z^2 - 1)^2$.

\originalnumber{3}{exercise:3.4.3.12.3}
3. Same problem for $w = z^3 - 3z$.


% Original PDF physical page 116
\index{Conformal mapping|)}\index{Fundamental region|)}\index{Riemann surface|)}\sourcepage{116}
